\chapter{Insurance service, finance and claims}
\label{ch:insurance}

An insurer receives cash before much of its service is delivered. Cash premium,
service revenue and profit therefore follow different schedules. A short
coverage example first connects gross insurance to reinsurance. The subsequent
two-year cases introduce unearned profit, loss allocation and finance one
change at a time.

\section{Insurance: cohorts and future service}

An insurer requires policy cohorts or groups with issue dates, coverage,
premiums, benefits, expenses, lapses, claims development, guarantees and
reinsurance relationships. Asset portfolios need their own contracts and
valuations. The liability state must distinguish expected future cash flows,
discounting, risk adjustments or margins, unearned profit measures where
applicable, incurred claims and movements between opening and closing states.
The IFRS 17 overview's distinction between measuring future cash flows and
recognising profit as service is provided motivates this separation
\citep{ifrs17}.

The module API accepts policy/cohort state, economic assumptions, non-economic
assumptions and information date. It returns a measurement vector, movement
attribution and proposed recognition events. Experience differences, changes
in assumptions, finance effects and service effects remain separately
identifiable. Reinsurance held has its own record and measurement; it is not
an automatic percentage offset to every insurance liability.

Scenarios can alter discount curves, mortality, morbidity, lapse rates,
inflation, catastrophe losses and policyholder behaviour. Their joint
dependence and cohort exposure must be declared. A portfolio total can match
history while hiding very different guarantee and duration profiles. This
makes alternative economic histories especially valuable, but also limits
what public-data reconstruction can identify. A full insurance package needs
actuarial and accounting review beyond the core release.

\clearpage
\section{Insurance service and reinsurance in a complete cohort}

The example uses one-year policies and a one-year proportionate treaty starting
when those policies are recognised, with uniform service for each. Both groups
qualify for premium allocation through the one-year conditions in paragraphs
53(b) and 69(b); paragraphs 62--62A determine treaty recognition. Assume no
acquisition cash flows, investment components, significant financing or onerous
remaining coverage. Claims and recoveries are expected within a year of the
claims being incurred. The inspected AASB~17 text incorporating IFRS~17 specifies
the movements in paragraphs 53--59, allocation in B126 and separate reinsurance
requirements in 63--70A \citep{aasb17}. Gross and ceded groups retain distinct
identities.

Begin with capital and cash of 20. Receive premium 120 and recognise a
remaining-coverage liability of 120. Pay reinsurance premium 24 and recognise
a coverage asset of 24. After half the uniform service, recognise revenue 60
by releasing gross coverage and reinsurance expense 12 by releasing its
coverage asset. Estimated incurred claims are 40, of which 30 is paid;
the remaining-claim risk adjustment is assumed to be 2. Gross claims expense
is therefore 42 and incurred-claims liability is 12. Assume the treaty
reimburses half the gross expected claims, recovery cash of 15 has arrived,
the risk transferred is valued at 1, and reinsurer non-performance has zero
expected effect in this example. The reinsurance claims asset is
$20+1-15=6$, and recovery income is 21.

Cash is $20+120-24-30+15=101$. Assets are cash 101, remaining reinsurance
coverage 12 and claims recovery 6, totalling 119. Liabilities are remaining
gross coverage 60 plus incurred claims 12, totalling 72. Profit is
$60-42-12+21=27$ and equity is $20+27=47$. These balances reconcile without
netting reinsurance against the insurance liability. If counterparty credit
risk changes, the recovery measurement changes through its own inputs; it
does not mechanically halve every gross liability or risk adjustment.

Longer-duration business adds a fulfilment-cash-flow vector, discount curve,
risk adjustment, contractual service margin, loss component and coverage-unit
schedule. Select portfolio and annual cohort, assess onerous grouping, project
cash flows within contract boundaries, discount with the declared financial
assumptions, measure non-financial risk separately, and attribute subsequent
movements to new business, experience, assumptions, finance and service.
Future-service changes adjust the supported service-margin operation subject
to the loss rules; current/past-service effects enter their specified result
components. Unsupported variable-fee, transition or loss-recovery operations
must stop those dependencies explicitly. A PAA example does not implement
these longer-duration methods by implication.

\enlargethispage{3\baselineskip}
An actuarial reserve adapter also needs a concrete estimation contract.
For cumulative claims $C_{i,k}$ by origin year $i$ and development year $k$,
the chain-ladder point method estimates each factor from matched observed
cells, $\widehat f_k=\sum_i C_{i,k+1}/\sum_i C_{i,k}$, then projects the
latest cumulative claim along the remaining factors. A triangular example
with rows $(100,150,180)$, $(120,180)$ and $(140)$ gives factors 1.5 and
1.2, outstanding amounts 36 and 112 for the unfinished rows, and total 148.
\citet[Sections 2--3, equations (1)--(3)]{mack1993} makes explicit the
conditional-mean and variance assumptions and independence across origin
years; its aggregate-error discussion retains dependence from shared estimated
factors. A single final-development observation cannot supply an empirical
tail variance without an additional assumption. The adapter therefore returns
the point projection and marks uncertainty estimation unavailable in this tiny
example. It supplies neither a full predictive distribution nor an IFRS risk
adjustment. Calendar effects, inflation, changing claims handling and tail
selection require separate sensitivity and validation before real use.
\hypertarget{group-state-and-general-measurement}{%
\section{Group state and general measurement}\label{method-insurance-methods-group-state-and-general-measurement}}

The following cases apply selected general-measurement and service-allocation provisions from the retained AASB 17 compilation \citep{aasb17}. They extend the premium-allocation example without presuming actuarial calibration or a complete insurer model.

The general method distinguishes fulfilment cash flows (FCF), including an
explicit risk adjustment (RA), from the contractual service margin (CSM),
which holds unearned profit. The liability for remaining coverage (LRC)
concerns future service; the liability for incurred claims (LIC) concerns
service and insured events already experienced. Keeping these components
separate makes a change in future estimates different from an unpaid claim.

Contract records retain insured risks, premiums/benefits, service and settlement dates, options/guarantees, cancellation/repricing rights and acquisition costs. Group admission evaluates portfolio/similar-risk management, onerous status and required issuance cohorts under AASB 17.14--24, recognition under 25--28 and boundary under 34. Cash outside the boundary belongs to future contracts, not this group's profit. GroupState retains expected dated cash flows, current discount curves, inception rates, explicit risk adjustment, CSM, loss component, LRC/LIC, coverage units, service experience and assumption vintages. Same premium totals are not sufficient state for merging different guarantees or claim tails.

General measurement follows 32--38: fulfilment cash flows are probability-weighted future net outflows, adjusted for timing/financial risk without double counting, plus explicit non-financial risk adjustment. CSM removes initial unearned gain under its conditions; an initial net loss is recognized under 47, not negative CSM. Paragraphs 40--44 separate remaining coverage and incurred claims. For a nonparticipating group, accrete CSM at the specified locked-in rate, apply future service changes at their required rate to CSM up to its available amount, record any excess loss and allocate current service over current plus future coverage units under B119. B96--B97 distinguish future service from incurred-claim changes, financial changes and current experience; not every revision is a CSM adjustment.

A MeasurementInput must resolve cash-flow/law, discount, risk-adjustment method, service grouping and coverage-unit evidence. Missing probabilities or RA support permit conditional sensitivity cases but block a source-supported measurement. Risk adjustment is the compensation the entity requires under 37; a standard deviation, confidence quantile or arbitrary percentage is only a candidate method until justified for that entity. Disclosures and the confidence-level equivalent, where applicable, are separate obligations. Market-consistent financial inputs and entity risk compensation must not be conflated.

\hypertarget{complete-two-year-general-measurement-case}{%
\section{Complete two-year general-measurement case}\label{method-insurance-methods-complete-two-year-general-measurement-case}}

Take a two-year nonparticipating group with one premium 120 received at inception, claims expected/actually incurred and paid 40 each year, zero discount rate, RA 10 initially released 5 with each equal service year, no acquisition costs, investment component, lapses, taxes, finance changes or reinsurance. Coverage units are 1 in each year because equal benefits and coverage are assumed. Before premium payment, expected net cash outflow is 80\ensuremath{-}120=\ensuremath{-}40, RA 10, FCF \ensuremath{-}30 and CSM 30, with net initial carrying amount zero. After the premium debit cash 120/ credit LRC 120. Its components are remaining expected claims 80, RA 10 and CSM 30.

Each year releases expected service claims 40, RA 5 and CSM 15: debit LRC 60/ credit insurance revenue 60. Incur claims 40: debit service expense 40/credit LIC 40; payment debits LIC 40/credits cash 40. End of year one: cash 80, LRC 60 (claims 40, RA 5, CSM 15), LIC 0 and equity/profit 20. End year two: cash 40, liabilities 0, cumulative equity/profit 40. Revenue 120 and claims 80 reconcile to premium 120 less claims cash 80. No investment component has been mistaken for revenue.

At year-one close, a separate pre-next-service change increases the remaining claim estimate from 40 to 60, RA unchanged. FCF rises 20; consume remaining CSM 15 and recognize service loss 5. LRC rises from 60 to 65, comprising FCF 65, CSM 0 and loss component 5 as a tracked portion, not an added liability on top. Debit service loss 5/credit LRC 5 with component transfer CSM\ensuremath{\rightarrow}FCF 15. A decrease of 8 in the revised future claim estimate, before any additional service/finance movement, first reverses loss 5 and restores CSM 3. FCF 57 plus CSM 3 gives LRC 60. Credit loss reversal 5/debit LRC 5. This tests 48--52 ordering. Later systematic allocation of loss component to releases/finance must reduce it to zero by coverage end; it cannot be amortized as revenue a second time.

In a separate case, an already-incurred unpaid claim estimate rises 40 to 48. Debit service expense 8/credit LIC 8; do not consume future-service CSM. A discount-rate change is financial under B97; its profit/OCI route follows the selected permitted finance policy. The zero-rate base case does not remove the need for separate current and locked-in rate calculations in a real group.

For the unreversed onerous alternative, carry the loss component through the remaining service year. Opening LRC 65 comprises future claims 60 and RA 5, with CSM 0 and loss component 5 within that LRC. A declared proportional allocation assigns 5/65 of each expected claim/RA release to the loss component, consistent with paragraphs 49--52 and B124's revenue exclusions. The expected claim release allocates 60/13 to loss reversal and the RA release 5/13; the total is 5, leaving loss component 0. Debit LRC 65, credit revenue 60 and credit loss reversal 5. Separately incur and pay claims 60 through LIC. Period profit is 60\ensuremath{-}60+5=5; combined with first-year profit 20 and the earlier onerous loss 5, lifetime profit is 20. Closing cash 20 equals equity 20 and all liabilities are zero. Treating all 65 as revenue while also reversing loss 5 overstates profit by 5.

\hypertarget{nonzero-discount-and-service-allocation}{%
\section{Nonzero discount and service allocation}\label{method-insurance-methods-nonzero-discount-and-service-allocation}}

A second complete case keeps premium 120, two equal coverage units and RA 10, but expected and realized claims are 44 after one year and 48.4 after two years. Both the inception and current discount curves remain a flat 10\% throughout; there are no new estimates, options, fees, taxes, investments or OCI election. This RA schedule is an authored measurement input, not an actuarial calibration: it declines by 5 per year, with its entire change in the service result under 81. All other insurance finance effects go to profit or loss. The retained source supports the mechanics in 44, 80--88, B119 and B123--B124.

At inception the present value of the two expected claims is
\[
 \frac{44}{1.1}+\frac{48.4}{1.1^2}=40+40=80.
\]
 FCF before premium is 80+10\ensuremath{-}120=\ensuremath{-}30, CSM 30, and after premium LRC 120. In year one the claim PV accretes 8; CSM accretes 3 at the locked-in rate. Allocate the resulting CSM 33 across the two current/future units, releasing 16.5. Revenue is expected claims 44 plus RA 5 plus CSM 16.5=65.5. Service expense 44 and finance expense 11 leave profit 10.5. After paying 44, cash 76 and LRC 65.5 (claim PV 44, RA 5, CSM 16.5) reconcile to equity 10.5.

In year two claim PV accretes 4.4 and CSM accretes 1.65. Release all remaining CSM 18.15, RA 5 and expected claims 48.4, giving revenue 71.55. Service expense 48.4 and finance expense 6.05 leave profit 17.1. Final cash/equity 27.6 equals premium 120 minus claims 92.4; all group liabilities are zero. Total revenue 137.05 less finance expense 17.05 equals premium 120. Setting revenue equal to cash premium each year or releasing 15 of CSM each year would ignore the specified financing. These are separate component calculations followed by debit finance expense/ credit LRC for accretion, debit LRC/credit revenue for service and the stated LIC recognition/settlement entries. Cash is deliberately assumed uninvested.

\section{Two issued groups and their reinsurance share one balance sheet}
\label{sec:insurance-connected-groups}

An insurer's accounts contain several groups at once. A profitable group
cannot lend its unearned margin to an onerous group merely because both
belong to the same company. Reinsurance may recover part of the loss, but
does so through a separately measured contract. The following extension
puts those three measurements on one balance sheet and follows their cash
through the two service years. It uses the retained AASB 17 rules for an
Australian private-sector insurer's calendar 2024 and 2025 periods.

Start with cash and equity of 100. Group A is the preceding nonzero-rate
example: premium 120, claims of 44 and 48.4 at the two year ends, initial
claim present value 80, risk adjustment 10 and CSM 30. Group B receives
premium 100 and expects claims of 66 and 72.6 at the same dates. At the
same 10\% rate, B's claim present value is
\(66/1.1+72.6/1.1^2=120\). Its separately supplied risk adjustment is 10,
so it recognises a loss of 30 and has no CSM. After premium receipt its
remaining-coverage liability is 130, including a loss component of 30.
Both groups contain contracts issued at the beginning of 2024 and provide
two equal annual coverage units. Their different profitability requires
separate groups under paragraphs 16 and 47; assume no pricing-law exception
under paragraph 20. Estimating cash flows together does not authorise pooling
their CSMs \citep[paragraphs 14--28, 32--38 and 47]{aasb17}.

The current and inception discount curves remain at the assumed 10\%.
Each issued group's risk adjustment releases 5 per year, entirely in the
service result under paragraph 81. All insurance finance enters profit or
loss under paragraph 88(a). Claims are incurred and settled at each year
end, with realised claims equal to the admitted expectations on this path.
Cash is uninvested. There are no acquisition costs, investment components,
taxes, currency effects, new contracts or estimate changes. The rates,
expected claims, service units and risk adjustments are supplied conditional
measurements; these calculations do not estimate them from insurer data.

Now add treaty R, entered into before B is recognised. It covers all B
claims at 50\%, with no exclusions, limits, settlement delay or expected
non-performance. Its premium of 70 is paid immediately. Expected recoveries
are 33 and 36.3, with present value 60. A separately supplied transferred-risk
adjustment of 5 releases 2.5 each year. Its equality to half B's adjustment
is an assumption of this example, not a general quota-share valuation rule.
Without the initial loss-recovery adjustment, the net cost held in R's CSM
would be \(70-60-5=5\), and the asset after payment would be 70.

Because R was entered into in time and covers all B's contracts, paragraphs
66A and B119C--B119D recognise recovery income of \(30\times0.5=15\).
The adjustment increases R's cost CSM to \(5+15=20\), so its complete
remaining-coverage asset after payment is \(60+5+20=85\). Its loss-recovery
component of 15 is already within that asset. It is not another receivable.
Debit B's initial service loss 30 and credit its liability 30; debit R's
asset 15 and credit recovery income 15. Receive gross premiums 220 and pay
the reinsurance premium 70. Cash is 250, issued liabilities are 120 and 130,
and the reinsurance asset is 85, leaving equity 85
\citep[paragraphs 60--68, 78 and B119C--B119F]{aasb17}.

\subsection*{Financing moves the loss component as well}

B's loss component must be allocated across claim releases, risk release
and finance under paragraphs 50--52. For this unchanged group, use the
initial proportion \(\alpha=30/130=3/13\) for all three. If \(F\) is
the period's financing of B's future claim value and \(S\) is its expected
claims plus risk-adjustment release, its loss component \(K\) follows
\[
 K_{\mathrm{close}}=K_{\mathrm{open}}+\alpha F-\alpha S.
\]
This is a declared systematic allocation for these facts. It preserves
\(K=\alpha\) times B's remaining-coverage liability and reaches zero at
coverage end. New estimates, coverage or risk patterns would require their
own reassessment of the allocation; the example does not prescribe this
ratio for every insurer \citep[paragraphs 49--52 and B124]{aasb17}.

In year one, B's claim value accrues finance of 12. Its loss component
therefore increases by \(36/13\), before service releases claims 66 and
risk adjustment 5. Of that release of 71, \(213/13\) reverses the loss
and \(710/13\) is insurance revenue. The closing liability is
\(130+12-71=71\), comprising future claim value 66 and risk adjustment 5.
The loss component is \(30+36/13-213/13=213/13\), within that liability.
Debit finance expense 12 and credit B's liability; then debit the liability
71, credit revenue \(710/13\) and credit loss reversal \(213/13\).
Separately incur and pay claims 66 through the incurred-claims liability.

R's claim value accrues 6 and its CSM accrues 2. The asset gains finance
income 8. Its CSM of 22 is then allocated across two equal current and
remaining service units, releasing 11. Recovery 33, risk release 2.5 and
CSM release 11 reduce remaining coverage by 46.5; the closing asset is
\(85+8-46.5=46.5\), comprising 33, 2.5 and 11. Its loss-recovery component
tracks half B's loss component: finance adds \(18/13\), service reverses
\(213/26\), and the closing component is \(213/26\). No component is
added on top of the measured asset. In this unchanged, fully covered case,
the component reversal is represented by the scheduled cash-flow and risk
release; paragraph 66(bb) must not produce a second CSM reversal for it.

The service entry debits reinsurance premium-allocation expense
\(46.5-213/26=498/13\) and loss-recovery reversal \(213/26\), and credits
R's remaining-coverage asset 46.5. Separately recognise recovery income 33
and an incurred-recovery asset 33, then collect the cash. Thus R's service
result is \(33-46.5=-13.5\), with finance income 8 shown separately.
The component determines presentation within that service result; it does
not change its total. Paragraph 86 also prevents the reinsurance premium
allocation from reducing issued insurance revenue
\citep[paragraphs 66--67, 80--88, B72, B119 and B119F]{aasb17}.

A retains its preceding result: revenue 65.5, claims expense 44 and finance
expense 11 give profit 10.5, with remaining coverage 65.5. B has service
profit 5 and finance expense 12, giving a further loss of 7. Adding R's
loss of 5.5 makes the year-one movement in equity \(-2\), after the
initial net loss of 15. Cash is \(250-44-66+33=173\), and
\(173+46.5-65.5-71=83\) confirms closing equity.

\subsection*{Settlement and timing expose different errors}

In year two A produces profit 17.1 as before. B incurs finance expense
6.6, then releases expected claims 72.6 and risk adjustment 5. Its loss
component first increases by \(99/65\), then reverses \(1164/65\), reaching
zero. Of the total release 77.6, revenue is \(776/13\) and the remainder
is loss reversal. Paying claims 72.6 leaves B with no liability and
period loss 1.6. R accrues finance income \(3.3+1.1=4.4\), then releases
recovery 36.3, risk adjustment 2.5 and all CSM 12.1. Its asset release
is 50.9; its loss-recovery component also reaches zero. Collecting 36.3
leaves reinsurance service expense 14.6 and net expense 10.2. Total
year-two profit is \(17.1-1.6-10.2=5.3\).

The following balances are after each date's premium or claim settlements.
All incurred claims and incurred recoveries have been settled; the exact
event sequence preserves those intermediate balances before the cash entries.
\begin{center}
\begin{tabular}{lrrr}
\toprule
 & Initial receipt/payment & Year one & Year two\\
\midrule
Cash & 250 & 173 & 88.3\\
A: remaining-coverage liability & 120 & 65.5 & 0\\
B: remaining-coverage liability & 130 & 71 & 0\\
R: remaining-coverage asset & 85 & 46.5 & 0\\
Equity & 85 & 83 & 88.3\\
\bottomrule
\end{tabular}
\end{center}
The final cash of 88.3 equals initial capital 100 plus gross premiums 220,
less reinsurance premium 70 and gross claims 231, plus recoveries 69.3.
Cumulative profit is \(-15-2+5.3=-11.7\). The two issued groups and the
held treaty finish with every measurement component at zero.

Suppose instead that R is entered into immediately after B's recognition,
with no intervening service, insured event or financing and otherwise the
same prospective coverage. It is a valid different history. It cannot
recover the original recognised loss under paragraph 66A because it fails
B119C's ordering condition. R then starts with cost CSM 5 and asset 70,
with no loss-recovery component. Opening equity is 70. Its year-one closing
asset is \(33+2.5+2.75=38.25\), giving total equity 74.75. Year-two cash
and equity still finish at 88.3. The difference is the timing of recognition,
not a difference in cash. Paragraph C20A’s transition exception does not apply to this ordinary
new-business example.

Incorrectly pooling A and B produces a third history of reported amounts.
Their combined premiums 220 equal claim present value 200 plus the two
supplied risk adjustments 20. A pooled calculation would therefore show
neither CSM nor initial loss, conceal B's loss of 30, and suppress the timely
treaty's recovery adjustment of 15. It reports equity 100 initially and
91.25 after year one, although final cash is again 88.3. Those accounts
balance but violate the separate-group recognition requirements. Likewise,
adding the recovery component to the full asset overstates opening equity
by 15, while omitting financing from loss allocation makes the component
run out before coverage ends. Intermediate group identity, component
movements and recognition time are essential checks; terminal cash alone
cannot distinguish these mistakes.

Actual use also requires the treaty's boundary, collectability, measurement
inputs and eligible contract mapping. Partial coverage needs the systematic
loss allocation in B119E; a nominal quota cannot supply it. Changes in those
facts or in the estimates require new measurement events under paragraphs
66--67, rather than extending this fixed schedule unchanged. The shared
event procedure commits journal and component movements for each dated
event together, while preserving separate issued and held group identities.
\citep[paragraphs 24, 63--67, 78--86 and B119C--B119F]{aasb17}


\section{Commit each insurance event as a complete book transition}
\label{sec:insurance-atomic}

Separate group measurements still have to produce one dated company balance
sheet. In the preceding example, R's loss-recovery component depends on B's
loss component. Publishing B's financing while retaining R's opening
measurement would combine amounts from different stages of the calculation.
The proposed close operation therefore declares A, B and R as the required
group set for each event in this example. Each group retains its own identity
and measurement; the event publishes their consequences together in one
entity, book and branch. This is a declared dependency for these facts, not a
requirement to update every group in every insurance transaction.

The year end contains four events, in order: finance, service, claim and
recovery incurrence, then settlement. Each has its own business-event
identifier and immutable opening state. Keeping those boundaries preserves
the incurred balances and the distinction between earning income and moving
cash. The following intermediate states follow from the preceding journals.
Coverage columns show liabilities for A and B and an asset for R.
\begin{center}
\begin{tabular}{lrrrrr}
\toprule
After event & Cash & A coverage & B coverage & R coverage & Equity\\
\midrule
Opening & 250 & 120 & 130 & 85 & 85\\
Finance & 250 & 131 & 142 & 93 & 70\\
Service & 250 & 65.5 & 71 & 46.5 & 160\\
Incurrence & 250 & 65.5 & 71 & 46.5 & 83\\
Settlement & 173 & 65.5 & 71 & 46.5 & 83\\
\bottomrule
\end{tabular}
\end{center}
Service releases the coverage amounts before the separate claim expenses
and recovery income are recognised. Incurrence creates liabilities of 44
and 66 and a recovery asset of 33. Those balances explain why equity is 83
at that stage even though cash remains 250. Settlement clears them and
moves cash by \(-44-66+33=-77\), leaving equity unchanged. Collapsing the
four events into a single final balance would conceal this distinction.

\subsection*{A valid component calculation is part of permission to post}

The proposed operation uses the common
\texttt{OperationRequest} for \texttt{accounting.commit@1}. Its four arguments
remain references to the accounting basis, posting authorization, journal
proposal and economic delta. The expected company head remains in the
common request envelope. Each group measurement identifies the group,
issued or held role, covered group where applicable, event phase, opening
and proposed closing components, book, branch, time, information view,
method and source. A complete measurement has actual component values;
an unresolved one has no closing value and gives its reason.

Before looking up a previous receipt, the operation validates the complete
group set, common ownership and information view, admitted evidence,
measurement calculation and posting authorization. If R is unresolved,
the whole declared event lacks permission. A and B's completed proposals
remain available for diagnosis, but they do not form a closing company
state. A missing group, duplicate group, wrong branch or different method
also prevents lookup. The existing receipt index cannot turn malformed
historical input into an eligible request. This proposed profile admits
only the authored synthetic case; submitting it to a real book fails that
permission check.

For an eligible new event, the journal and economic delta must then agree
with the measurements. During finance, R's asset rises by 8 because its
claim present value gains 6 and its cost CSM gains 2. Its loss-recovery
component separately rises by \(18/13\), as a component within the asset.
An extra debit to R's asset and credit to finance income for \(18/13\)
would balance, yet overstate both the asset and income. The operation
rejects that proposal. It also rejects a proposal whose journal is correct
but whose economic delta leaves the recovery component at its old value.
The monetary accounts and component memory must both be correct.

Publication inserts the journal, complete component and account changes,
closing company state, their dependency records and edges, and the immutable
receipt, then advances the company head in the same transaction. The finite
exchange examples list those published records separately from the global
registry of proposed records. Merely retaining a rejected measurement for
inspection therefore does not make it part of the book. A rejected request
publishes none of its proposed transition. This requirement has one legal
book as its owner; it supplies no transaction guarantee across different
books or entities.

\subsection*{Retry the event without repeating its accounting}

The business-event lookup key remains the entity, book, branch and
business-event identifier. After input validation, an equal request
fingerprint returns the original receipt before testing the freshness of
today's head. Suppose all four events have committed, advancing the head
from generation zero to generation four. Retrying the original finance
request returns its original receipt and closing state at generation one.
The response separately reports the observed generation-four head and that
this invocation made no additional commit. It does not apply finance again
or replace today's state with the old closing state.

Equality covers the complete request, including its request identifier,
budget, original expected head and dependency identities. Changing any of
those bytes while retaining the business-event key is a conflict, not an
exact retry. A previously unseen event with an old head is stale. The
transaction repeats both the key and head checks before publication: an
identical request that has just won returns its receipt; different request
bytes for that event conflict; a different event that advances the head
makes this request stale. In each case the competing transaction's changes
are distinguished from any changes made by the requester.

If the acknowledgement disappears, the client cannot infer whether the
transaction committed. The exchange therefore records a transport
observation with unknown commit status, the original submission and its
exact request bytes. It does not invent a received operation result. One
possible storage state still has the opening head and no receipt; the other
has the complete closing state and receipt. Resubmitting the exact request
commits once in the first case and returns the old receipt in the second.
A timeout supplies no reason to change the event identifier, opening head
or economic facts.

These finite examples check the equations, complete record bindings and
required decisions for both storage states. They specify observable
behavior; actual transaction isolation, crash recovery and durability need
later implementation tests. Their completed coverage is a decision on one
request, with no interpretation as probability mass over financial histories.
The underlying source applicability and supplied actuarial assumptions of
the connected-group example continue to govern its accounting meaning.


\hypertarget{paa-reinsurance-and-valuation-methods}{%
\section{PAA, reinsurance and valuation methods}\label{method-insurance-methods-paa-reinsurance-and-valuation-methods}}

The preceding premium-allocation case states separate eligibility and recognition timing for gross insurance and reinsurance. PAA under 53--59 requires the one-year coverage or nonmaterial-difference condition; a short payment lag is not the eligibility test. LRC starts with received premium and applicable acquisition adjustments, then releases insurance revenue under the prescribed service pattern. LIC still measures incurred obligations using the applicable fulfilment-cash-flow rules; the allowed short-settlement discount simplification needs its own conditions. Emerging onerous facts require a loss assessment; a premium balance cannot be presumed sufficient.

A one-year policy can leave claims unpaid for several years. The distinction
between the period of insurance service and the period until settlement is
therefore essential even when PAA is available. The official ASBE~25 text
separates three timing tests. Article~50 permits PAA when each contract's
coverage is no longer than one year, or when the stated no-material-difference
condition holds. Article~54 concerns a different interval: the time between
each portion of service and the related premium due date. If that interval
does not exceed a year, the remaining-coverage liability need not be adjusted
for financing. Article~56 instead measures from claim incurrence to expected
payment or receipt. Its option to omit the time-value and financial-risk
adjustment requires settlement within a year of incurrence
\citep[Articles 50, 54, 56]{mofasbe25}. The retained AASB~17 compilation makes
the corresponding distinctions in paragraphs 53, 56 and 59(b).

Consider a separate group of issued policies with exactly one year of coverage.
All amounts are illustrative currency units. Premium 120 is due and received
when coverage starts, with no earlier onerous trigger; group recognition is
therefore at that date under Article~16. Assume valid insurance grouping,
uniform risk release, no acquisition costs, investment component, reinsurance,
tax, other expenses or other movements, and no facts indicating onerous
remaining coverage. Elect the permitted remaining-coverage financing exception.
After half a year, recognise revenue 60 and retain remaining coverage of 60.
Articles~53 and 57 support this roll-forward and time allocation; Article~57
requires a different allocation when risk release is not predominantly
time-based \citep{mofasbe25}. These assumptions isolate the claim measurement
from other changes in the group's result.

At that reporting instant a claim is first incurred, with an expected payment
of 60.5 exactly two years later. This is a mean estimate, not a certain
settlement. Assume an appropriate flat annual effective discount rate of
10\% and a separately estimated non-financial risk adjustment of 5 in current
measurement units. The inputs are authored for the example; applying them to
an insurer would require market and actuarial evidence. Articles~21, 23, 25
and 26 require current probability-weighted cash flows, appropriate discounting
and a separate non-financial risk adjustment \citep{mofasbe25}. The claim's
cash-flow present value is $60.5/1.1^2=50$, so its liability is $50+5=55$.
The risk adjustment is not discounted a second time. The claim has only just
been incurred, so no intervening finance accrual enters this measurement.

Record debit cash 120/credit remaining coverage 120 on premium receipt, then
debit remaining coverage 60/credit insurance revenue 60 for the service
provided. Recognising the claim gives debit insurance service expense 55/credit
incurred claims 55, consistent with Article~31. The group then has cash 120,
remaining coverage 60, incurred claims 55 and a cumulative result of 5:
$120-60-55=60-55=5$. These are balances attributable to the group at mid-year,
not a complete insurer forecast. Using the undiscounted $60.5+5=65.5$ for this
two-year claim would overstate the required liability by 10.5 and understate
the result by the same amount. The shortcut is wrong under the stated rules
and assumptions. A zero rate would remove that discounting difference while
leaving the separate risk adjustment intact. Articles~86(3) and 97 also require
disclosure of the distinct claim components and the PAA/financing choices
\citep{mofasbe25}.

This comparison uses identified editions. The Ministry page reproduces the
2020 ASBE~25 text; Article~120 states 1 January 2023, but that sentence does
not settle every entity group's adoption schedule.
Section~\ref{sec:insurance-adoption} distinguishes the Ministry's later
explanation of that timetable from a company's reported adoption. The retained
December 2022 AASB~17 compilation applies to annual periods beginning on or
after 1 January 2023 but before 1 July 2026. The shared result in this example
establishes a scoped measurement comparison. Applying it to a particular entity
still requires its applicable rules and contract facts; later amendments and
broader equivalence require separate evidence.

\subsection{Acquisition costs and claim finance follow different clocks}
\label{sec:insurance-paa-cost-finance}

The one-year coverage example establishes PAA eligibility without shortening
the claim's settlement period. Acquisition costs add another choice. Under
the historical IFRS~17 text, costs incorporated into the PAA liability reduce
its opening balance and are subsequently amortised as an expense. Immediate
expense is permitted under 59(a) when each contract's coverage period at
initial recognition is no more than one year
\citep[IFRS~17.53--59]{eu20231803}.

Consider a separate eligible one-year group with premium 120 and acquisition
cash paid 12 at recognition. Opening cash and equity are each 200. Assume
uniform service and risk release, no investment component, tax or FX, and
the permitted LRC financing exception because the service-to-premium-due
interval is at most a year. No claims are incurred before coverage ends,
and the base case has no onerous indication. Cash after receipt and costs
is 308. Deferring the acquisition cost gives opening LRC \(120-12=108\)
and equity 200. Expensing it gives LRC 120 and equity 188.

At mid-year, both cases recognise revenue 60. In the deferred case, debit
acquisition expense 6 and credit LRC 6 for amortisation, then debit LRC 60
and credit revenue 60. LRC is \(108+6-60=54\), and cumulative profit
is 54. In the expensed case, LRC is 60 and cumulative profit is
\(60-12=48\). After the second half's service, both LRC balances are
zero and both cases have recognised revenue 120 and total acquisition
expense 12. Equal final costs do not justify erasing the interim choice.
Time allocation also requires the supplied uniform risk pattern; B126--127
instead uses expected incurred-service-expense timing when the risk-release
pattern differs significantly \citep[IFRS~17.55,59(a),B126--127]{eu20231803}.

An adverse mid-year branch requires a fresh onerous calculation. If current
FCF for the remaining coverage is 70 while deferred-cost LRC is 54, recognise
loss 16 and increase LRC to 70. Adding the entire FCF 70 to the existing
54 would overstate the liability by 54. The measurement must use the
discount treatment required by paragraph~57, including its connection to
the LIC simplification \citep[IFRS~17.57--58]{eu20231803}.

Return to the base case. At coverage end a claim is incurred, with expected
payment 48.4 two years later, current discount rate 10\% and separately
supplied risk adjustment 2. Initial LIC is \(48.4/1.1^2+2=42\);
cash 308 less LIC 42 leaves equity 266. The claim-incurrence rate of 10\%
differs from the group's assumed inception rate of 5\%. If the portfolio
elects to split LIC finance between profit and OCI, B133 uses the
claim-incurrence rate specified in B72(e)(iii), not the inception rate
\citep[IFRS~17.59(b),B72(e)(iii),B133]{eu20231803}.

One year after incurrence, suppose the current rate falls to zero and stays
there. The claim's current present value becomes 48.4, while a separate
schedule using the claim-incurrence rate reaches 44. Finance expense is
8.4 in total: 4 in profit and 4.4 in OCI. Assume the risk adjustment
releases 1 in each of the two years, wholly in the service result under
paragraph~81. First-year LIC is then 49.4 and equity 258.6. In the second
year, current claim value remains 48.4, so total finance is zero; the
claim-incurrence schedule reaches 48.4 through profit expense 4.4 and OCI
income 4.4. Release the remaining risk adjustment and pay the claim.
Final cash and equity are 259.6, equal to opening 200 plus premium 120
less acquisition costs 12 and settlement 48.4. Cumulative OCI is zero.
Using the inception rate would give first-year profit finance of only 2
and the wrong allocation, even if the total liability were correct.


\subsection{Reinsurance and direct participation}

Reinsurance has its own group, boundary, recognition time, counterparty risk, cash-flow, CSM/cost and possibly loss-recovery component under AASB~17 paragraphs 60--70A. Do not multiply every gross insurance liability by a quota share: excluded claims, limits, deductibles, timing, reinstatements and counterparty non-performance matter. Claim 100 with an eligible 50\% recovery but expected 90\% collectability gives conditional expected recovery 45, not 50; altering transferred insurance risk and reinsurer non-performance are different inputs. This simple calculation does not establish treaty recognition or RA by itself. Gross liability 100 remains; reinsurance recovery is a separate asset/measurement with its own authority.

For direct participation, test B101's conditions using the actual underlying items, substantial-share and variation judgments. Retain underlying fair values, entity fee share and guarantee/risk-mitigation choices; paragraphs 45 and B101--B118 specify the variable-fee adjustments to CSM, with onerous and mitigation exceptions. A local component illustration of a qualifying share: underlying-value increase 20 and entity's fixed 10\% share gives candidate variable-fee change 2; eligibility, guarantees and other FCF changes must be evaluated before posting. The full option/guarantee valuation and transition calculations are further supported profiles, with source and numerical obligations, not an assumed linear fee model.

Eligibility depends on the contract and the insurer's expectations at inception.
The policyholder must participate in a clearly identified pool, with the
required substantial share of fair-value returns and substantial variability
in payments. The variability assessment covers the contract's duration on a
probability-weighted present-value basis. A minimum guarantee can bind in some
scenarios without automatically excluding direct participation. Later market
movements do not reopen the classification; contract modification follows its
own rules. The identified reference pool also need not be held by the insurer,
so the contractual link and the actual asset holdings require separate records
\citep[paragraphs B101--B108 and 72]{aasb17}.

\clearpage
A small settlement example shows why the guarantee matters even when the
stated share is fixed. Consider only a maturity-benefit component, in one
currency at one date. Assume a held pool allocated to this component has amount
$U$, and the policyholder is owed the larger of 90 and $0.9U$. These are
illustrative contractual choices. The amount left from the pool after payment
is $U-\max(90,0.9U)$. Writing the payment as
$0.9U+\max(90-0.9U,0)$ shows that the amount left is
$0.1U-\max(90-0.9U,0)$: the nominal retained share less the extra payment
needed to meet the minimum.

\begin{center}
\begin{tabular}{@{}rrr@{}}
\toprule
Pool amount & Policyholder payment & Amount left\\
\midrule
80 & 90 & $-10$\\
100 & 90 & 10\\
120 & 108 & 12\\
\bottomrule
\end{tabular}
\end{center}

Moving from 80 to 100 leaves 20 more after payment: the ordinary share rises
by 2 and the extra floor payment falls by 18. Moving from 100 to 120 leaves
only 2 more because the minimum no longer binds. The negative amount at 80
is a funding shortfall for this component. These unweighted settlement
outcomes explain the exposure; they do not establish insurance-contract
eligibility, a present value or the insurer's profit. Other benefits and costs
remain outside this component example, and its held-pool assumption does not
replace the standard's separate accounting for investment assets
\citep[paragraphs B66(a), B107--B108]{aasb17}.

The accounting calculation therefore needs more than this payoff schedule.
Paragraph B104 defines the variable fee as the entity's share of underlying
fair value less fulfilment cash flows that do not vary based on underlying
returns. Changes in the obligation to pass through the underlying fair value
do not adjust CSM, whereas the entity-share changes relate to future service.
Financial-guarantee changes have their own future-service treatment under
B113(b). Paragraph 45 retains the loss-component and risk-mitigation exceptions,
as well as new-contract, currency and service-allocation movements. Applying
those rules requires the group's measured cash flows, discount and risk
adjustments, remaining CSM and loss component, service information and any
eligible mitigation designation. The original increase of 2 remains a valid
share-component calculation; the complete CSM movement needs these additional
inputs \citep[paragraphs 32--38, 45, 48--50 and B104--B118]{aasb17}.

\section{Pricing the floor before releasing the margin}
\label{sec:insurance-priced-guarantee}

The settlement schedule identifies a liability that increases when the
underlying assets fall. To measure that liability before settlement, we need
a financial model. A small, explicitly specified market makes the distinction
between a price and a forecast visible. The retained standard permits a
replicating portfolio to measure cash flows that it matches in amount, timing
and uncertainty; other techniques must remain consistent with observable
market prices \citep[paragraphs B44--B48]{aasb17}.

Suppose a non-dividend-paying traded asset costs $S_0=100$. Its next observed
value is either $S_u=120$ or $S_d=80$, after which it remains at that value
until payment. Cash earns zero interest, so its accumulation factor is $R=1$.
There are exactly two financial states, unrestricted borrowing and shorting,
no transaction costs and no counterparty failure. These are assumptions of
the illustration. The financial observation precedes all supplied coverage
units; the subsequent service interval contains the mortality risk. An actual
period with service already provided would require its own dated allocation.

For payoffs $X_u,X_d$, let $\delta$ denote the units of the traded asset and
$b$ the amount deposited in cash. Matching the two payments requires
\[
 \delta S_u+bR=X_u,\qquad \delta S_d+bR=X_d.
\]
Subtracting the second equation from the first gives
$\delta=(X_u-X_d)/(S_u-S_d)$. Substitution then determines $b$ and the
portfolio's current cost:
\[
 b=\frac{S_uX_d-S_dX_u}{(S_u-S_d)R},\qquad
 V_0=\delta S_0+b
     =\frac{qX_u+(1-q)X_d}{R},\qquad
 q=\frac{RS_0-S_d}{S_u-S_d}.
\]
Thus the weight $q$ follows from today's asset and cash prices. Positive
state prices require $S_d<RS_0<S_u$, with $R>0$. Here $q=1/2$.
Additional financial states or restrictions that prevent the matching trades
would invalidate this particular valuation method.

The policyholder's investment payment is $\max(0.9S,90)$. Its top-up above
the ordinary share is $(90-0.9S)^+$, where $x^+=\max(x,0)$.
That top-up pays zero in the up state and 18 in the down state.
The matching portfolio has $\delta=-0.45$ and $b=54$:
$-0.45(120)+54=0$ and $-0.45(80)+54=18$.
Its cost is $-0.45(100)+54=9$. Adding the ordinary share's value of 90
gives a financial-benefit value of 99. This derives the option value; no
option price has been supplied as an unexplained input.

For a separate physical forecast, suppose the up state has probability
$3/4$. The expected top-up is then $18/4=4.5$, although its market-consistent
price remains 9. Suppose also that a covered death triggers an additional
payment of 20, with physical probability $1/10$, independently of the market
state. Its expected cost is 2. Nonfinancial risk adjustment of 1 is a
separately supplied measurement assumption, with its entire release in the
service result. The financial state prices, physical joint law and compensation
for nonfinancial risk therefore occupy different records
\citep[paragraphs 33, 37, 81 and B42--B49]{aasb17}.

\subsection{The same measurement through service and payment}

Use the retained historical Australian provisions for a conditional
calendar-2024/25 example. Assume the contract has significant insurance risk,
an enforceable link to the identified pool and the required direct-participation
characteristics. The component rights terminate together; neither can be
retained separately after the other matures. This supplies the specified
interrelationship under B32(b), together with an assessment that no embedded
derivative must be separated. These admission findings need actual terms
before real use; neither $q$ nor the 90\% share proves them
\citep[paragraphs 10--13, B17--B23, B31--B32 and B101--B108]{aasb17}.

The investment payment is due whether the insured lives or dies. There are
no surrender, lapse, acquisition or tax cash flows in this example, and all
figures are pretax. The insurer holds one unit of the underlying asset at
fair value through profit or loss and makes no risk-mitigation election.
The portfolio used to calculate the guarantee price is not its actual hedge.
Insurance finance also enters profit or loss under paragraph 89(a).
The investment component remains within the insurance measurement but is
excluded from service revenue and expense under paragraph 85
\citep[paragraphs 85, 89 and B115--B118]{aasb17}.

Start with owner capital and cash of 50, and receive premium 110.
Before receipt, fulfilment cash flows are $99+2+1-110=-8$, giving CSM 8.
After receipt, the remaining-coverage liability is 110: financial benefits 99,
mortality estimate 2, risk adjustment 1 and CSM 8. Buying the underlying
for 100 leaves cash 60 and investment 100 against liability 110 and equity 50.

At the financial observation, the future-service fee change is the change
in the insurer's 10\% share less the change in the top-up value
\citep[paragraph 45 and B111--B114]{aasb17}. In the up state this is
$2-(0-9)=11$, increasing CSM from 8 to 19. Fulfilment cash flows become
$108+2+1=111$ and the liability becomes 130. Investment gain 20 and
insurance finance expense 20 offset.

In the down state the fee change is $-2-(18-9)=-11$. The existing CSM
absorbs 8; the excess 3 becomes a service loss. Fulfilment cash flows and the
liability are $90+2+1=93$, with CSM zero and a loss component of 3
\emph{within} that liability. Investment loss 20 and insurance finance income
20 offset, while the service loss reduces equity to 47. The loss component
does not create a second liability, and the binding floor does not reopen
the inception VFA classification
\citep[paragraphs 45, 48--52, 87(c) and B102]{aasb17}.

Completing the service interval releases the expected mortality cost, risk
adjustment and remaining CSM. Up-state revenue is $2+1+19=22$.
In the down state the service-related release of 3 reverses the loss component,
leaving revenue zero. Actual mortality expense is either zero or 20, first
recognised through the incurred-claim liability and then paid. Selling the
underlying and paying the investment component of 108 or 90 completes all
balances.

\begin{center}
\begin{tabular}{@{}lrrrr@{}}
\toprule
Joint outcome & Physical weight & Profit & Final cash/equity & OCI\\
\midrule
Up, survival   & $27/40$ & 22  & 72 & 0\\
Up, death      & $3/40$  & 2   & 52 & 0\\
Down, survival & $9/40$  & 0   & 50 & 0\\
Down, death    & $1/40$  & $-20$ & 30 & 0\\
\bottomrule
\end{tabular}
\end{center}

Each row is a complete conditional history. Operating cash is premium 110
less the investment benefit and mortality payment; the investment purchase
and sale give investing cash of $+20$ or $-20$. This classification assumes
nontrading investments under the retained cash-flow provisions; fair-value
measurement alone does not determine trading purpose
\citep[paragraphs 14--16]{aasb107dated}. The physical expected profit is 14.5
under the supplied law, but the four rows preserve the separate downside
exposures. An unknown mortality law, risk adjustment, dependence assumption
or admission finding leaves a named unresolved measurement rather than a
default probability or zero cost.

The exact state prices, journals, component changes and adverse alternatives
are retained with the acquisition/treaty cases in
\path{data/accounting-master/}\allowbreak\path{finite-method-decisions-20261015/insurance/}.
The companion method, \emph{Pricing the guarantee and retaining separate
insurance groups}, binds valuation, probability and risk-adjustment records to the same group,
date and evidence view, while keeping their meanings distinct. General market
models and actuarial calibration require additional methods and validation.



\section{The variable fee separates investment returns from service}
\label{sec:insurance-variable-fee}

The preceding guarantee valuation measures a contractual payoff. Direct
participation adds a service question: which changes belong to the insurer's
unearned fee? Admission requires all three B101 conditions, assessed using
expectations at inception: an enforceable share of a clearly identified
pool, a substantial share of fair-value returns payable to policyholders,
and substantial variability of their payments with those returns. The
assessment considers the contract's duration and probability-weighted
present values, including scenarios in which a guarantee binds. It is
reassessed only upon a relevant modification. Reinsurance issued or held
cannot qualify as direct participation under B109
\citep[IFRS~17.B101--109]{eu20231803}.

For an admitted group, the source separates the obligation equal to the
underlying items' fair value from the variable fee deducted for future
service. Changes in the former do not adjust CSM. Changes in the insurer's
share and in eligible future-service costs do, subject to loss and risk
mitigation conditions. The latter adjustments use current rates; a GMM
locked-rate calculation cannot be copied into this model
\citep[IFRS~17.45,B104,B110--114]{eu20231803}.

Suppose a qualifying insurance group receives premium 100, invested in its
identified pool. Policyholders are entitled to 80\% of the pool's value
regardless of whether an insured event occurs, together with separate
insured benefits. Supply remaining nonfinancial claims of 10 and risk
adjustment 2, two equal current/future service units, and zero discount,
tax and FX effects. The 80\% amount is an investment component under the
stated terms. All eligibility and service-quantity findings are explicit
case premises. With pool value 100, FCF is \(80+10+2=92\), leaving
CSM 8 and LRC 100. Opening cash and equity 200 become cash 200, pool
asset 100 and liability 100 after receiving and investing the premium.

Before the first service allocation, the pool rises to 120. The
policyholder's share rises by 16 to 96. The insurer's 20\% share rises
by 4, increasing CSM to 12. FCF becomes \(96+10+2=108\), and total
LRC becomes 120. The financial increase in the liability is therefore 20,
not merely the policyholder's increase of 16. With all pool gains and
insurance finance in profit, asset gain 20 offsets insurance finance
expense 20; the extra fee 4 remains unearned within CSM.

Now supply expected and actual first-period insured claims 5 and risk
adjustment release 1. Allocate the CSM 12 equally across the current and
remaining service units, releasing 6. Revenue is \(5+1+6=12\), claims
expense 5 and service profit 7. Paying the claims leaves cash 195, pool
asset 120 and LRC \(96+5+1+6=108\), giving equity 207. In the second
period, with unchanged pool value and the same claims and release, profit
is another 7. Sell the pool for 120 and repay the investment component 96.
Final cash and equity are 214. The repayment reduces the liability and
cash; it is excluded from both insurance revenue and service expense
\citep[IFRS~17.81,83--85,B119]{eu20231803}.

\subsection{Losses and risk mitigation change the CSM adjustment}

Consider a separate adverse branch from the initial FCF 92 and CSM 8,
before any service. Future nonfinancial costs rise by 15. Use the CSM 8
first, then recognise service loss 7 and a tracked loss component 7 within
LRC 107. Equity falls from 200 to 193. A subsequent cost reduction of 6
reverses loss 6 and leaves loss component 1, CSM zero and LRC 101. A
further reduction of 4 reverses the last loss 1 and restores CSM 3;
FCF 97 plus CSM 3 brings LRC back to 100. Those loss reversals remain
service results under the specified rules even when a qualifying source
of the original loss was a decline in the insurer's underlying-item share
\citep[IFRS~17.45,48--50,87(c)]{eu20231803}.

A financial guarantee change needs a different cause record. From the same
initial position, suppose its supplied future-service value increases by
6 while a derivative gains 6. Without applying the risk-mitigation choice,
FCF rises to 98 and CSM falls to 2, leaving LRC 100. The derivative gain
enters profit, increasing equity to 206. If B115--116's conditions are met
and the insurer applies the choice to this change, CSM remains 8 and LRC
rises to 106. Insurance finance expense 6 then offsets the derivative gain,
leaving equity 200. The choice changes recognition timing, not the
derivative's value or the obligation.

This treatment requires a previously documented risk-management objective
and strategy, an economic offset and credit risk that does not dominate
that offset. B115 specifies which derivatives, FVTPL non-derivative
instruments or held reinsurance can mitigate each exposure. It cannot be
elected retrospectively because a mismatch became inconvenient. B117A
takes the relevant finance to profit for derivative/FVTPL mitigation and
specifies the matched policy for held reinsurance. When B116's conditions
cease, B118 stops the treatment prospectively without undoing previous
recognition \citep[IFRS~17.B115--118]{eu20231803}.


\hypertarget{partial-reinsurance-and-changes-in-insurance-estimates}{%
\section{Partial reinsurance and changes in insurance estimates}\label{partial-reinsurance-and-changes-in-insurance-estimates}}

The following methods use IFRS 17 in the retained September 2023 EU compilation. They extend the preceding nonparticipating general-measurement examples. Contract rights, group assignments, recovery percentages, risk adjustments and discount rates are explicit conditional inputs. A change in one of these inputs must retain its cause: changed service expectations, financial risk, counterparty performance and a new legal agreement do not have interchangeable accounting consequences.

\hypertarget{only-the-covered-contracts-support-loss-recovery}{%
\subsection{Only the covered contracts support loss recovery}\label{only-the-covered-contracts-support-loss-recovery}}

IFRS~17.62--66B \citep[IFRS~17.62--68 and B119C--F]{eu20231803} keeps a held treaty separate from the issued contracts, with its own entry, coverage and recognition dates. Proportionate coverage can delay recognition until an underlying contract is recognised. B119C permits the initial onerous-loss recovery only if the treaty was entered into before or at the same time as the underlying onerous contracts were recognised. The record must distinguish these dates; the treaty's later cash payment does not identify when it was entered into.

When a treaty covers only some contracts within an onerous issued group, B119E requires a systematic and rational allocation of the group's loss to the covered contracts. B119D then multiplies that covered loss by the percentage of their claims expected to be recovered. A treaty's headline quota percentage must not be applied to uncovered losses. Preserve the contract-to-issued-group and contract-to-treaty joins, covered risks and periods, attachment points, limits and exclusions, the loss-allocation method, expected claims/recoveries and the separate transferred-risk adjustment. For excess-of-loss or capped cover, expected recoveries require the admitted claim distribution and treaty payoff; a constant quota share cannot be used without support.

Take an issued group with premium 100, expected claims 118 and risk adjustment 2, all measured with zero discounting at recognition. Initial loss is 20 and CSM is zero. A source-supported allocation in this synthetic contract set identifies a covered portion with premium 41, claims 48 and risk adjustment 1, giving loss 8; the uncovered portion has premium 59, claims 70 and risk adjustment 1, giving loss 12. These amounts sum to the group; they are not separate groups created to improve the result.

A treaty entered into on time covers half of the covered contracts' claims. Expected recovery is 24, transferred-risk adjustment is a separately supplied 0.5, and premium paid is 26.5. After payment, the ordinary fulfilment asset is 24.5 and cost CSM is 2. Loss recovery is

\[8\times 50\%=4,\]

so cost CSM rises to 6 and the measured held asset to 30.5. The loss-recovery component of 4 sits inside that asset; it is not another receivable. Starting from cash and equity 50, receive issued premiums 100 and pay treaty premium 26.5. Cash 123.5 plus held asset 30.5 equals issued liability 120 plus equity 34. Applying 50\% to the entire loss would recognise 10 and overstate equity by 6. If this same treaty had been entered into only after the underlying recognition, the initial loss-recovery adjustment would not apply.

\hypertarget{a-changed-loss-estimate-affects-both-groups-without-netting-them}{%
\subsection{A changed loss estimate affects both groups without netting them}\label{a-changed-loss-estimate-affects-both-groups-without-netting-them}}

Before any service or time passes, reduce expected claims on the covered contracts from 48 to 44, leaving other assumptions unchanged. Issued fulfilment liability falls from 120 to 116. Paragraph 50(b) first reduces the loss component, from 20 to 16, and recognises a loss reversal of 4; there is no new issued CSM while that loss component remains.

Expected held recovery falls from 24 to 22. Under paragraph 66(c)(i) \citep[IFRS~17.66(c),66(bb),B119F]{eu20231803}, this change follows the underlying change that did not adjust the issued CSM, so it reduces the held asset and creates expense 2 rather than being hidden in held CSM. Held fulfilment asset becomes 22.5, cost CSM stays 6 and total held asset becomes 28.5. The covered loss is now 4, so B119F's recoverable loss-component limit is 2 and the loss-recovery component falls from 4 to 2. That component change tracks the same recognised recovery loss; posting a second expense for it would double count the event. Equity becomes \(123.5+28.5-116=36\), an increase of \(4-2=2\).

If all remaining services are then provided and claims equal these revised estimates, issued liability 116 is released through insurance revenue 100 and loss reversal 16; claims expense and payment are 114. Issued service-period profit is 2. Held coverage asset 28.5 is released as premium allocation 26.5 and reversal of loss recovery 2, with incurred recovery 22 and its collection recorded separately. Held service-period loss is 6.5. The entity ends with cash and equity \(123.5-114+22=31.5\) and no group balances. Initial combined loss 16, estimate-change gain 2 and service-period loss 4.5 sum to lifetime loss 18.5, equal to the change from opening equity 50. Reinsurance has not removed or netted away the issued obligation.

B119F requires later loss-recovery components to reflect changes in the underlying loss component and never exceed the portion expected to be recovered. Track separately adjustments already included in fulfilment-cash-flow changes and reversals under 66(bb) that are not. The cap does not authorise inventing recovery income. Under paragraph 67, a change in reinsurer non-performance risk does not adjust CSM as future service; preserve collateral and dispute assumptions required by paragraph 63. Transferred-risk adjustment must be measured under paragraph 64, not mechanically copied as a fraction of the issued adjustment.

\hypertarget{changed-terms-and-changed-estimates-are-different-events}{%
\subsection{Changed terms and changed estimates are different events}\label{changed-terms-and-changed-estimates-are-different-events}}

For a treaty change, retain the old and new legal terms, effective date, covered-contract links and consideration. An exercise of a right already contained in the contract is not a modification under paragraph 72. A new treaty requires its own recognition and grouping decision and cannot be backdated to create initial-loss recovery. For a modification, apply paragraph 72's tests \citep[IFRS~17.72--77]{eu20231803} to the actual changed contract: changed scope or separated components, a substantially different boundary, a different initial group, changed direct-participation status, or loss of eligibility for a previously used premium-allocation approach can require derecognition and a new contract. If none applies, paragraph 73 treats the resulting cash-flow changes as estimate changes, with the held-contract modifications in paragraphs 63--68 where applicable.

For required derecognition, paragraphs 76--77 remove the affected cash flows and risk adjustment, adjust the remaining group's CSM under the applicable model and revise remaining coverage units. A transfer or qualifying modification uses paragraph 77's actual transfer consideration or equivalent new-contract premium adjustment rather than simply deleting the old balance. The equivalent premium is a required valuation input, not the old carrying amount chosen for convenience. Missing terms or that valuation leaves the particular modification unresolved. The generic tests identify the required procedure; they do not price an unspecified changed treaty.

\subsection{The replacement premium connects the old and new groups}
\label{sec:insurance-modification-values}

A qualifying modification changes both the removed obligation and the new
contract. Paragraph~77 prevents the insurer from choosing the new opening
balance simply by copying the old one. For a contract derecognised under
paragraph~72, determine the premium that would be charged for an equivalent
new contract at the modification date, then subtract any additional premium
actually charged for the modification. This difference adjusts the old
group's CSM together with the removed fulfilment cash flows, and is treated
as a premium received when measuring the new contract
\citep[IFRS~17.72--77]{eu20231803}.

Take a general-measurement group with FCF 100, comprising expected claims
90 and risk adjustment 10, and CSM 30. Cash is 200, so liability 130 leaves
equity 70. A legally effective modification substantially changes one
contract's boundary and requires a new contract under paragraph~72. Its old
allocated FCF is 40, comprising claims 36 and risk adjustment 4. Assume zero
discount rates and no intervening service, tax, FX or OCI. The independently
supplied equivalent new premium is 60, of which 10 is actually charged and
received as an additional premium. The new obligation has FCF 52, comprising
claims 48 and risk adjustment 4. These values are case inputs; they are not
inferred from the old carrying amount.

The deemed premium is \(60-10=50\). Removing FCF 40 while transferring
value 50 reduces the old group's CSM by 10. Its remaining FCF is 60 and
CSM 20, giving liability 80. The new contract's total consideration of 60
less FCF 52 gives CSM 8. One accounting sequence debits the old liability
50 and credits the new liability 50, then debits cash 10 and credits the
new liability 10 for the actual receipt. Cash 210 equals old liability 80,
new liability 60 and unchanged equity 70. Recording a second cash receipt
for the deemed 50 would invent liquidity.

If the old group previously had three equal remaining coverage units and
the removed contract contributed one, it now has two; the new group has its
own service assessment. A proportional removal of 40\% of the old CSM
would leave 18 instead of 20 and would disregard the equivalent-premium
rule. Over the remaining service, assume actual claims equal the supplied
estimates and all remaining risk adjustment and CSM are released. The old
group produces profit \(80-54=26\); the new group produces
\(60-48=12\). Final cash \(210-54-48=108\) equals opening equity 70
plus profit 38.

The same removed FCF with an actual third-party transfer premium of 50
instead reduces the old liability by 50 and pays cash 50; there is no new
issued contract in this separate case. Cash 150, liability 80 and equity
70 reconcile. If the original CSM were only 6, either adjustment requiring
10 would consume it and recognise service loss 4. For the modification,
old liability 60 plus new liability 60 would then leave equity 90 from
cash 210, compared with opening equity \(200-106=94\). The resulting
loss component 4 is tracked within the old liability, not added again.
Paragraphs~44(c), 48--50 and 77 govern that limit
\citep[IFRS~17.44,48--50,77]{eu20231803}.

Where the removed contract carries allocated accumulated OCI, paragraph~91
adds a separate presentation decision. Amounts arising under the systematic
allocation in 88(b) are reclassified to profit on the specified transfer or
derecognition; amounts arising under the held-underlying-item policy in
89(b) are not. Preserve that origin and the allocated amount before closing
the old contract. The no-OCI example above does not eliminate this step
\citep[IFRS~17.91]{eu20231803}.


\hypertarget{future-service-consumes-csm-before-creating-a-loss}{%
\subsection{Future service consumes CSM before creating a loss}\label{future-service-consumes-csm-before-creating-a-loss}}

For issued contracts without direct participation, paragraph 44 and B96--B97 \citep[IFRS~17.44,48--52,B72,B96--97]{eu20231803} distinguish future nonfinancial service changes from financial changes and incurred-claim experience. Use initial-recognition discount rates for the specified CSM adjustments under B72(c), while fulfilment cash flows use current rates. Changes in current rates, financial risk and incurred-claim estimates do not enter CSM by simply labelling them future cash flows. Preserve the current and initial-rate valuations and the reason for each change; their difference is part of the finance analysis, not an unrecorded balancing amount.

To isolate the nonfinancial rule, assume no passage of time, zero rates, no new contracts and no service release. A group has fulfilment liability 80, CSM 20 and total liability 100. An adverse future-service estimate of 15 raises fulfilment liability to 95 and reduces CSM to 5, with no immediate profit effect. A further adverse change of 10 raises fulfilment liability to 105, exhausts CSM and creates loss 5 and a loss component of 5. A subsequent favourable change of 8 first reverses that loss 5; only the remaining 3 restores CSM. Closing fulfilment liability 97 plus CSM 3 is 100. Crediting all 8 to profit would overstate the reversal by 3. The loss component is an internal allocation of the liability, never an additional liability of 5.

\hypertarget{finance-presentation-preserves-total-income-and-expense}{%
\subsection{Finance presentation preserves total income and expense}\label{finance-presentation-preserves-total-income-and-expense}}

For the nonparticipating case in paragraphs 87--90 \citep[IFRS~17.87--91 and B128--132]{eu20231803}, a portfolio-level policy can present insurance finance entirely in profit or disaggregate it using the required systematic allocation. Under B129--B131, the choice and the contract characteristics determine the method. If financial assumptions do not substantially affect the nominal amounts paid, the systematic allocation uses the initial-recognition rates specified by B72(e)(i). OCI is the difference between that allocation and total finance, not a discretionary smoothing reserve. Over the group's duration it totals zero under B130.

Consider a GMM group with a now-incurred, fixed claim of 121 payable in two years, no remaining coverage or CSM, and zero risk adjustment because no nonfinancial uncertainty remains in this illustrative obligation. Initial and current discount rates are 10\%, so liability is \(121/1.1^2=100\). The next year's current rate falls to zero; one year remains, and current liability becomes 121. Total first-year finance expense is 21. Under the admitted OCI policy, the initial-rate allocation is 10 in profit and the difference 11 in OCI. In the second year, with rate zero unchanged, current liability stays 121, so total finance is zero. The initial-rate allocation accretes the alternative carrying amount 110 to 121, giving profit expense 11 and OCI income 11. Paying 121 extinguishes the liability. Cumulative profit finance is 21 and OCI is zero.

With opening cash 150, liability 100 and equity 50, first-year equity falls to 29 and stays 29 in the second year because profit expense and OCI income offset. Final cash and equity are 29. Presenting all finance in profit gives expenses 21 and zero instead, with the same total equity and cash. The OCI choice changes presentation, not the contractual claim or current measurement. This is the B131 fixed-payment case; variable-payment B132, VFA/underlying-item B134--B136, PAA and actual curve/RA/process valuation retain their separate methods. A successful fixed-claim example does not establish those models.


\section{Variable benefits require a different finance allocation}
\label{sec:insurance-variable-finance}

The fixed-claim example uses initial rates because financial assumptions
do not substantially change the nominal payment. B132 covers a different
case: financial assumptions substantially affect the amounts paid. For the
specified OCI policy, the cash-flow finance allocation may use a rate that
spreads revised expected finance over the remaining duration at a constant
rate. Contracts using a crediting rate have the alternative based on
amounts credited and expected to be credited. Nonparticipating CSM finance
retains its prescribed initial-rate calculation; any separately allocated
risk-adjustment finance must follow the source's consistent allocation
\citep[IFRS~17.B130--132]{eu20231803}.

Take a group with no remaining service, CSM or nonfinancial uncertainty,
but a payment that remains sensitive to financial assumptions. Its liability
and the amount measured using the systematic allocation both start at 100.
A new financial assumption raises the expected payment in two years from
121 to 144. Supply a market-consistent current valuation of
\(144/1.1^2=14400/121\), approximately 119.0083. Under the chosen
constant-rate allocation, the remaining rate solves
\(100(1+r)^2=144\), so \(r=20\%\). The remeasurement at that instant
produces OCI expense approximately 19.0083; no time has yet elapsed for
the revised allocation to accrue interest.

With the current discount rate staying at 10\%, current liability rises
to \(144/1.1\), approximately 130.9091, after one year. Total finance
expense is approximately 11.9008. The separate allocation rises from 100
to 120, giving profit expense 20 and OCI income approximately 8.0992.
In the final year, current liability rises to 144, giving total finance
approximately 13.0909; the allocation rises from 120 to 144, giving profit
expense 24 and OCI income approximately 10.9091. All amounts are checked
as exact fractions before display rounding. Total profit finance is 44,
and accumulated OCI is zero. Opening cash 200 and equity 100 close at
cash and equity 56 after payment. Retaining the original 10\% allocation
would reach only 121 and leave an unexplained 23 at settlement.

For a general payment schedule, retain the revised dated cash flows,
carrying amount under the allocation, remaining duration and solved rate,
then reconcile the full expected allocation and OCI unwind. A numerical
failure or non-unique admissible solution remains unresolved; it is not
permission to use an arbitrary rate. Current measurement and systematic
allocation are separate calculations even when their values happen to agree.

\subsection{Held underlying items determine the matching presentation}

When a direct-participation insurer holds the underlying items, paragraph
89(b) permits a different OCI policy. B134 places in profit insurance
finance that exactly offsets the profit income or expense on those items,
with both amounts presented separately. The residual insurance finance
goes to OCI. Revisit the preceding pool increase of 20, but now supply a
valid asset measurement that reports 8 in profit and 12 in OCI. The matching
insurance finance is expense 8 in profit and expense 12 in OCI, eliminating
both net financial movements. The insurer's fee increase still enters CSM
under the service method; matching is not an additional fee adjustment
\citep[IFRS~17.87A--90,B134]{eu20231803}.

Actual holdings must be recorded separately from the pool named by the
contract. If holdings change the available choice between 88(b) and 89(b),
B135--136 preserves the previously accumulated OCI under a separate
reclassification schedule based on the former approach and assumptions at
the change date. Those assumptions are then frozen for that schedule.
Prior comparatives are not restated and accumulated OCI is not recomputed
as if the new policy had always applied. New-period finance follows the
newly applicable policy. A method that simply replaces the old rate and
reserve with a newly computed one would discard required history
\citep[IFRS~17.B135--136]{eu20231803}.

\section{A claim process needs joint probabilities}
\label{sec:insurance-joint-process}

The preceding calculations take cash-flow and risk-adjustment values as
inputs. Obtaining them requires a law for the dated obligations, including
their dependence. Means alone can support a linear expected-claim total
while giving the wrong answer for a treaty with a deductible. IFRS~17
requires probability-weighted discounted cash flows, with supportable
non-market assumptions and consistency with market variables. It does not
make independent claim draws a default when their dependence is unknown
\citep[IFRS~17.B37--41,B49--55]{eu20231803}.

Let claims \(X_1\) and \(X_2\) at two specified dates each be either 0
or 40, with probability one half for either amount. Assume zero discount
and no counterparty loss solely to isolate dependence. Total expected
claims are always 40. A treaty pays \(\max(X_1+X_2-60,0)\). If both
claims always occur together, total claims are either 0 or 80 and expected
recovery is 10. If exactly one claim occurs, total claims are always 40
and recovery is zero. The two models have identical single-date
distributions.

All joint laws on this declared binary support can be described by
\(a=\Pr(X_1=40,X_2=40)\). The two marginal constraints force the
following probabilities:
\[
\begin{array}{c|cc}
 & X_2=0 & X_2=40\\ \hline
 X_1=0 & a & 1/2-a\\
 X_1=40 & 1/2-a & a
\end{array}
\qquad 0\leq a\leq 1/2.
\]
Nonnegativity gives the displayed interval. Only the final cell produces
a recovery, of 20, so its expectation is \(20a\), ranging exactly from
0 to 10 within this finite family. Independent claims correspond to
\(a=1/4\) and recovery 5; independence is one member, not a conclusion
from the marginals. Substituting mean claims into the treaty would give
\(\max(40-60,0)=0\), which is wrong for every \(a>0\).

The same law supplies a sequential event procedure. Draw the first claim
with probability one half. After observing 40, the conditional probability
that the second claim is 40 is \(2a\); after observing zero it is
\(1-2a\). These transitions reproduce the joint table and provide a
conditional recovery forecast after the first event. This is an exact
finite-law specification, not a simulated result. If the evidence leaves
\(a\) unresolved, retain its admissible interval and separate conditional
results; do not assign a probability distribution over possible values of
\(a\) without additional evidence.

A general valuation method must similarly specify exposure and contract
boundaries, event times and transitions, observation conditioning, claim
severity, lapse and expense dependence, payment delays, discount treatment
and financial-variable dependence. Its version must distinguish changes
in the law from new observations under the same law. Risk adjustment then
requires an entity-supported compensation method and diversification
assumptions; the mean, variance or stop-loss value alone does not identify
that compensation \citep[IFRS~17.B65,B74--78,B86--89]{eu20231803}.
The exact family above supplies a small test of such a method's dependence
handling. Selecting and calibrating a richer law still requires the target
population, credible data and downstream validation specified in the
empirical protocols.


\section{Claims development estimates the remaining payments}
\label{sec:claims-development-estimation}

The claims-development example above supplies an estimation input, not a
recognition rule. Its point projection is 148, but the single observed final
factor leaves its variance unresolved. In the general method the factors are
\[
 \widehat f_k=\frac{\sum_{i\in O_k} C_{i,k+1}}{\sum_{i\in O_k}C_{i,k}},
\]
where $O_k$ contains origin years with both required observed cells. Each
incomplete row is projected through the remaining factors. The same triangle
therefore has ultimate claims 180, 216 and 252, with unpaid amounts 0, 36 and
112. Tail development, calendar inflation, zero or negative cells and dependence
need explicit extensions. The conditional means, variance scaling and
independent-origin-year assumptions behind the inspected Mack error formulas
\citep[equations (1)--(3), Theorems 1--3]{mack1993} cannot turn this tiny
triangle into a predictive law or a justified insurance risk adjustment.

\section{A reserve estimate is not a future payment}
\label{sec:reserving-bootstrap}

The reserve of 148 leaves an important question unanswered: how much might
the insurer ultimately pay? Repeating an estimation procedure can show how
its fitted answer varies. Even if that answer were known exactly, however,
the future claims would still vary. England's retained 2001 working
addendum makes this distinction operational: first resample and refit to
obtain future conditional means, then simulate payments conditional on
those means. The inspected text is Actuarial Research Paper 138; its exact
equivalence to the separately published 2002 journal addendum has not been
established \citep[Sections 2--3 and Appendix A]{england2001addendum}.

Consider a small probability example, separate from the triangle. Suppose
a fitted future mean, denoted by $M$, is either 2 or 8 with equal
probability. These are authored possibilities, not estimates from an
insurer. If $M=2$, let a payment $Y$ be 0 or 4, equally likely;
if $M=8$, let it be 4 or 12, equally likely. Within each fitted case,
the payment has mean $M$ and variance $2M$. Across cases, the fitted mean
has expectation 5 and variance 9. The total variation in payments includes
both the variation left within a case and the variation between case means:
\[
 \operatorname{Var}(Y)
 =\E\!\left[\operatorname{Var}(Y\mid M)\right]
  +\operatorname{Var}\!\left(\E[Y\mid M]\right)
 =\E[2M]+\operatorname{Var}(M)=10+9=19.
\]
One can check the result directly: $Y$ is 0, 4 or 12 with probabilities
$1/4$, $1/2$ and $1/4$. Its mean is 5, and its variance is
$(25/4)+(1/2)+(49/4)=19$. Reporting the distribution of fitted means
alone would instead give variance 9 and no payment above 8. The full
authored payment law gives probability $1/4$ to that event.

Even those two moments do not select a unique payment law. Keep the same
two fitted means, but, conditional on $M=2$, pay 1 or 6 with probabilities
$4/5$ and $1/5$; conditional on $M=8$, pay 6 or 16 with those probabilities.
The conditional means remain 2 and 8. Their variances are
$(4/5)1+(1/5)16=4$ and $(4/5)4+(1/5)64=16$, respectively.
The overall mean and variance remain 5 and 19, but the chance of paying
more than 8 is now $1/10$. Matching a reserve and a prediction error
therefore cannot determine a tail probability or the insurer's compensation
for non-financial risk.

\subsection{What the two stages require}

The addendum starts from a particular model for incremental payments.
For an observed incremental payment $X_{ij}$ in origin period $i$
and development period $j$, its fitted mean $m_{ij}$ is strictly positive.
The model assumes $\E[X_{ij}]=m_{ij}$,
$\operatorname{Var}(X_{ij})=\phi m_{ij}$ and
$\log m_{ij}=c+\alpha_i+\beta_j$, with reference levels fixed to identify
the parameters. This incremental model and Mack's conditional cumulative
model have separate assumptions, even where their fitted reserve amounts
agree. A common point estimate does not make their uncertainty models
interchangeable \citep[equations 2.1--2.3]{england2001addendum}.

The first stage measures each observed increment's departure from its fitted
mean using the Pearson residual $r_{ij}=(X_{ij}-m_{ij})/\sqrt{m_{ij}}$.
With $N$ observed cells and $p$ fitted parameters, the paper estimates
$\widehat\phi=\sum r_{ij}^2/(N-p)$ and resamples residuals after multiplying
them by $\sqrt{N/(N-p)}$. The adjustment must be applied once. It cannot
be followed by another correction of the same estimation variance
\citep[equations 2.4--2.6 and 3.1]{england2001addendum}.

Start from the observed latest cumulative amounts, work backwards through
fitted development factors, and difference the fitted cumulative amounts
into increments. Calculate the residuals and dispersion from those increments.
For each bootstrap draw, resample adjusted residuals with replacement and
invert the residual relation to form pseudo-increments. Accumulate them,
refit the development factors, project future cumulative claims and difference
them into future means. Retain the same fitted draw across all cells. This
produces a joint array of conditional means, not yet a joint array of payments
\citep[Section 2 and Appendix A, Triangles 1--9]{england2001addendum}.

The second stage chooses a process law. For a positive future mean $m$
and positive dispersion $\phi$, the paper proposes $\phi Z$, where $Z$
has a Poisson law with mean $m/\phi$. Its mean is $m$ and variance is
$\phi m$. A gamma law with shape $m/\phi$ and scale $\phi$ has the
same two moments but different support and tail. That is a modeling choice,
not a change of random-number syntax. Each future payment is drawn
conditionally, and aligned payments are summed by origin, payment period
and total \citep[Section 3 and Appendix A, Triangles 9--10]{england2001addendum}.

Dependence must survive aggregation. In the first small example, suppose
two cells share the same $M$ but are independent conditional on it. Their
conditional total variance is $4M$, with expectation 20, while their
conditional total mean is $2M$, with variance 36. Their total variance
is therefore 56. Independently redrawing each cell from its marginal
payment law would give $19+19=38$, losing the dependence created by the
shared fit. Marginal distributions are insufficient instructions for a
joint simulation.

\subsection{Sparse data, failed draws and implementation choices}

The six observed increments in the earlier triangle fit the incremental
model exactly. It has $p=2(3)-1=5$ parameters and only $N-p=1$ residual
degree of freedom, giving zero fitted dispersion. This arithmetic does
not establish that future payments are deterministic. It also does not
identify Mack's final development variance. The proposed interface returns
\code{DISPERSION\_UNRESOLVED} for this illustrative input and retains the
point reserve of 148 separately. More data or a justified additional
assumption is required before a predictive distribution can be admitted.

Every draw records its model, source version, data vintage, residual pool,
fitted parameters, process law and random stream. Nonpositive degrees of
freedom, an unsupported negative mean, a zero projection denominator or
non-finite arithmetic receives an explicit failure. Structurally absent
triangle cells have a separate mask. A failed payment is not converted
to zero or silently replaced by another draw. In the authored example,
if the $M=8$ case cannot be evaluated, its probability $1/2$ remains
unresolved: the known contribution to the full mean is $\tfrac12(2)=1$,
and a complete mean is withheld. In a finite bootstrap, attempted
failure frequency is reported as such; it does not identify population
probability.

The retained source itself needs a qualification. Its negative-mean
suggestion draws $X$ with mean $|m|$ and then subtracts $2m$. Read
literally for $m<0$, the resulting mean is $|m|-2m=-3m$, not $m$;
at $m=-4$ it is 12. That prescription is wrong relative to its stated
mean target in the inspected working version. A translation or reflection
can preserve a negative mean, but changes the support and requires a
different variance model. AccDSL leaves that extension unsupported
pending an explicit signed-claims model; this local finding does not
reject the positive-mean two-stage procedure
\citep[Section 3, final paragraph]{england2001addendum}.

Software names also conceal consequential choices. In the inspected
ChainLadder 0.2.17 source snapshot, \code{od.pois} calls a negative-binomial
generator for dispersion greater than one. The paper's scaled Poisson
and that negative binomial can share a mean and variance while assigning
different probabilities to possible payments. The same snapshot uses
signed gamma draws and replaces missing process values by zero. An adapter
must distinguish structural absence from failed computation, identify the
exact law and its parameter domain, and preserve separate resampling and
process streams before its output can satisfy this proposal. These are
source-code observations, not completed runtime validation
\citep{chainladder2023snapshot}.

The paper compares 1,000 simulations with analytic and earlier bootstrap
summaries. That provides a worked implementation comparison under its
assumptions. It does not establish future tail coverage, select an insurer's
risk adjustment or justify 1,000 draws as a universal budget. A later
empirical protocol must address calendar effects, tail development, pooling,
heldout diagonals, simulation uncertainty and consequences for measured
insurance groups \citep[Section 4, Tables 1--3 and Figure 1]{england2001addendum}.

\subsection{The measurement date changes the question}

A lifetime payment distribution differs from a one-year change in
estimated liabilities. The 2019 accepted manuscript by England, Verrall
and W\"uthrich defines the claims development result as opening reserves
less payments during the year and less the revised closing reserve. The
closing estimate uses the information that has emerged during that year
\citep[equations 5--7 and Section 3.2]{england2019lifetime}.
If the opening reserve is 100, payments are 30 and the revised closing
reserve is 80, the result is $100-30-80=-10$. Subtracting cash alone
would yield 70 and answer a different question. The scenario must carry
both realized next-year payments and the prescribed re-estimation rule,
with the correct information cutoff. A lifetime draw cannot be relabeled
as a one-year result. Neither statistic by itself supplies the discounting,
group allocation or risk compensation required for insurance measurement.

\section{What a bootstrap must reproduce}
\label{sec:bootstrap-consistency}

Separating fitted means from future payments is necessary, but it still leaves
the statistical target to be chosen. An insurer knows an observed claims
triangle. A method assessment may instead hold only its latest diagonal fixed,
allowing different earlier histories to lead to those same cumulative amounts.
These are different conditioning questions. Let $D$ denote the observed
triangle, $Q$ its latest diagonal, $R$ the future reserve payment and
$\widehat R(D)$ the estimate calculated from the triangle. A distribution of
$R-\widehat R(D)$ conditional on $Q$ allows the earlier entries of $D$, and
therefore the estimate, to vary. Conditional on the whole $D$, that particular
estimate is fixed.

A small authored example makes the distinction concrete. Suppose two earlier
histories are equally likely given the same $Q$ and produce estimates of 90
and 110. Let future payments be 80 or 120 with equal probability, independent
of those histories conditional on $Q$. The prediction error given $Q$ is then
$-30,-10,10,30$, each with probability $1/4$: its mean is zero and variance
500. Fixing the history whose estimate is 90 instead gives errors $-10,30$,
mean 10 and variance 400. Both calculations are correct for their stated
conditioning information. They cannot serve as interchangeable validation
targets. Nor does holding an estimate fixed imply that all uncertainty about
model parameters has disappeared; a Bayesian parameter distribution, for
example, would be a further explicitly specified model.

Steinmetz and Jentsch analyze a Mack bootstrap against the prediction error
conditional on the latest diagonal. Their asymptotic construction adds older
accident years and additional development observations while retaining the
current calendar boundary. It does not obtain more information merely by
increasing the number of bootstrap draws. The framework assumes positive
initial cumulative claims with finite moments, specified independence,
development factors tending to one, sufficiently fast decay of development
variances, and further moment and regularity conditions. It is distinct from
the incremental overdispersed-Poisson procedure in the preceding section
\citep[Sections 2--4, especially Assumptions 2.2--2.5]{steinmetz2024bootstrap}.

Their analysis separates an unscaled process term from a smaller estimation
term. Under the stated assumptions, the process term has a nondegenerate,
generally non-Gaussian limit. Correctly reproducing its distribution requires
the correct conditional family for future development factors; matching two
moments within a wrongly chosen family is insufficient. The original
bootstrap can reproduce that limit while missing the conditional mean and
covariance of the estimation term at its own, square-root sample-size scale
\citep[Assumptions 4.6--4.9, Theorems 4.8--4.11 and Remark 4.12]{steinmetz2024bootstrap}.

There is no contradiction in these two findings. For illustration, let a
process term be an equally likely $-1$ or $1$, and an independent estimation
term an equally likely $-1/n$ or $1/n$. Their sum has variance $1+1/n^2$.
Dropping the estimation term produces the same limiting distribution for the
unscaled sum, yet multiplying that omitted term by $n$ would have left variance
one. This example explains the distinction between scales; it is not a
derivation of the paper's Mack limit or evidence that a finite insurer triangle
is adequately approximated by it.

\subsection{Reconstructing the past requires a backward law}

The paper's proposed repair holds $Q$ fixed, simulates earlier cumulative
claims backwards, refits the resulting triangle, and centers simulated future
reserve payments at that refitted estimate. Future development uses the
originally estimated forward factors. The purpose of reconstructing the past
is to retain the effect of conditioning on the observed diagonal when
measuring estimation uncertainty. Theorem 6.3 requires Assumption 6.2,
including conditional limit properties for the backward reconstruction.
Appendix C invokes those properties; it does not establish them for every
positive distribution chosen by a programmer
\citep[Section 5, Assumption 6.2 and Appendix C]{steinmetz2024bootstrap}.

Forward conditional moments do not determine backward conditional moments.
Let $X$ be an earlier cumulative claim and $Y$ the next cumulative claim.
Suppose $X=1$ or $2$ with equal probability. Given $X=1$, let $Y=1$ or $3$,
equally likely. Given $X=2$, let $Y=3$ or $6$ with probabilities $2/3$ and
$1/3$. These positive claims satisfy
$\E[Y\mid X]=2X$ and $\operatorname{Var}(Y\mid X)=X$.
Thus the forward factor $Y/X$ has conditional mean 2 and variance $1/X$,
the form of a one-step Mack moment relation.

Now observe $Y=3$. The joint masses for $X=1$ and $X=2$ are $1/4$ and
$1/3$, respectively. Their sum is $7/12$, so Bayes' rule gives
$\Pr(X=1\mid Y=3)=3/7$ and $\Pr(X=2\mid Y=3)=4/7$.
The backward conditional mean is $11/7$, and its variance is $12/49$.
Dividing the observed cumulative amount by the forward mean factor would
instead give $3/2$. Choosing a backward factor with mean $1/2$ and variance
$1/3$ would imply a variance of 3 for $X$, also wrong for this joint law.
The example establishes the missing implication; it does not assert that this
finite distribution satisfies every assumption of an infinite-triangle theorem.

In the published paper, Section 5.1 proposes reciprocal-factor backward
moments, while Remark 5.1 and Assumption 6.2 separately require adequate
backward conditional behavior. The proposal therefore keeps such a generator
as a hypothesis requiring its own justification. Given a joint model, the
backward distribution must be derived from that model or approximated with a
checked error criterion. Choosing a gamma family and matching reciprocal
moments alone cannot establish the required conditional limit properties
\citep[Section 5.1, Remark 5.1 and Assumption 6.2]{steinmetz2024bootstrap}.

\subsection{The center and the comparison must use the same quantity}

Reserve payments exclude claims already paid. Write $q_i$ for the latest
cumulative payment in origin period $i$, $U_i^+$ for its simulated ultimate
payment, and $\widehat U_i^+$ for the ultimate payment predicted from a
reconstructed past. The reserve error must subtract like quantities:
\[
 R^+-\widehat R^+
 =\sum_i(U_i^+-q_i)-\sum_i(\widehat U_i^+-q_i)
 =\sum_i(U_i^+-\widehat U_i^+).
\]
For $q=100$, simulated ultimate payment 130 and fitted ultimate payment 120,
the reserve is 30, its estimate is 20, and their difference is 10.
Subtracting the ultimate estimate 120 directly from the reserve 30 gives
$-90$, shifting the error by the already-paid amount.

This distinction also resolves a literal inconsistency in the retained
published text. Equation (5.3), printed on page 96, writes the centering
quantity as a sum of projected ultimates without subtracting the latest
diagonal, although the preceding step defines reserves net of that diagonal.
Read literally, those two formulas produce the wrong reserve error. The
adjacent product-difference formula and decomposition (6.1) instead cancel
the diagonal and agree with the derivation above. The proposed implementation
uses that explicitly derived reserve center; it does not copy (5.3) literally.
No claim about the authors' unavailable execution code follows
\citep[Section 5.1, equation 5.3 and equation 6.1]{steinmetz2024bootstrap}.

The final journal article also changes the preprint's simulation comparison.
Section 7.1.1 generates whole triangles to assess unconditional prediction
interval coverage. Section 7.1.2 assesses conditional coverage using an
approximate backward generator in equation (7.4), based on the same
reciprocal-moment idea. The article acknowledges that the exact backward
conditional law is generally difficult to obtain. Its conditional simulation
therefore evaluates that constructed experiment; it does not by itself verify
coverage under the exact backward law implied by every forward model.
The published numerical differences remain reported simulation observations,
without a local replication or statistically supported ranking here
\citep[Section 7.1, Table 3 and Figures 1--3]{steinmetz2024bootstrap}.

For AccDSL, a reserving-method specification must identify the payment
horizon, conditioning information, process family, parameter-estimation
procedure, backward law when one is needed, and reserve-centering convention.
Validation must distinguish whole-triangle repeated-sample coverage,
coverage conditional on a diagonal, and prediction conditional on a full
observed history. A simulator with a convenient reverse rule is a candidate
model, not an independent reference for that same rule. Exact centering and
support checks precede simulations; target-specific comparison, uncertainty
estimation and downstream financial-statement checks precede any promotion.
The small triangle remains an accounting and arithmetic illustration, with
its unresolved dispersion preserved.

\section{A zero payment and a failed draw need different records}
\label{sec:reserving-adapter}

The distinction between payments and fitted means must survive the software
boundary. Suppose a reserving calculation requests two draws, each containing
payments for three future cells. In an authored example, the joint payment
vectors are $(0,4,8)$ and $(2,6,10)$. Their totals are 12 and 18, and the
sample mean is 15. These are possible nonnegative integer payments chosen
to illustrate the exchange; they are neither fitted to the preceding triangle
nor generated by a locally executed model.

Now let the last cell of the first draw fail. That draw's total is unresolved,
while the second total remains 18. Reporting 18 as the requested sample mean
would discard the failed draw. Replacing the failed value by zero would instead
produce totals 4 and 18, and mean 11. Neither number is the mean of the
requested complete sample. A diagnostic may state that the one completed
draw has mean 18, provided it retains the failed draw and withholds the full
sample mean. The fraction $1/2$ describes completed computation, not a known
probability mass assigned to unsuccessful economic histories.

This failure is relevant to the retained \code{BootChainLadder} source. Before
forming its returned process triangles and totals, it replaces every process
\code{NA} by zero. The returned array has no corresponding mask distinguishing
failed future draws from structural cells. A valid payment vector $(0,4,0)$
and a vector $(0,4,\text{failed})$ therefore become the same numeric vector
and total. Inspecting the final array cannot recover the lost distinction
\citep[\code{BootstrapReserve.R}, process-triangle coercion and returned fields]{chainladder2023snapshot}.

The unmodified output is consequently retained as source evidence but is
unsupported as a complete predictive result under this proposal. A future
instrumented implementation could preserve statuses before coercion, but
would need its own source identity, dependency lock and numerical checks.
Merely wrapping the returned array is insufficient. The documentary fixtures
specify that boundary; they do not establish an implemented or tested R adapter.

\subsection{The exchange retains the law and the whole draw}

An optional \code{reserving.evaluate@1} signature uses the common operation
envelope. Its declared profile names the claims triangle, method and process
law, requested draw identifiers, information view and producer evidence.
The present profile is a specification example that consumes authored
evidence. It makes no financial-book mutation and grants no probability-law
admission. Its request and result references use the same typed identity and
dependency rules as the other operations.

The triangle carries the currency, base payment unit, scale, origin and
development labels, payment-period mapping, valuation date and explicit cell
mask. Observed cumulative amounts, future cells and unavailable observed
amounts stay distinct. The inspected routine rejects unequal row and column
counts, so the proposed initial adapter subset is square even though the
native error message describes a broader shape. Rectangular support requires
a separately inspected implementation.

Each replicate carries one fit identifier and its complete joint vector of
future cells. A cell retains both its fitted mean and its process payment,
with a status and reason if the payment is unavailable. A valid zero remains
a value. Missing cells, repeated draw identifiers, a different fit inside one
draw, or a currency mismatch prevent acceptance. Any failed cell leaves the
draw total unresolved. Native evidence lacking the original status mask cannot
be relabeled as complete by an adapter.

The process-law identity includes its formula. With positive mean $m$ and
dispersion $\phi>1$, the inspected \code{od.pois} helper calls a negative
binomial with mean $m$ and size $m/(\phi-1)$. At $\phi=1$ it calls a
Poisson distribution. For $0<\phi<1$, its negative-binomial size would be
negative. A gamma method with positive shape and scale is a different
possibility; the helper's limitation does not reject that separate law
\citep[\code{rpois.od} and the process-law branches]{chainladder2023snapshot}.

Nor can the negative-binomial path silently satisfy a request for the
scaled-Poisson law from England's working addendum. For $m=2$ and
$\phi=3/2$, both have mean 2 and variance 3. The scaled-Poisson support
is $0,3/2,3,\ldots$, whereas the negative binomial can pay 1.
The base payment unit thus matters to the model, not just to presentation.
The positive-law profile excludes the package's separate signed handling;
a different signed-claims model would need its own support and validation.

\subsection{Stopping preserves what was actually computed}

A result reports requested, attempted, completed, failed and pending draws
separately. Its coverage domain is the set of computational replicate IDs.
It is not the company's family of economic histories. Complete arrays do
not establish calibrated forecasts, and an incomplete set of arrays does
not justify a probability distribution conditional on computational success.

If cancellation follows the first complete draw in the example, retain its
total 12 and mark the second draw unattempted. The diagnostic completed-draw
mean is 12; the requested sample mean remains unresolved. A stop request alone
does not establish that the worker has stopped. Unknown termination is
reported as partial with unresolved worker state. Resume requires a durable
checkpoint that binds the semantic manifest, original draw IDs and random
generator state; the native function's seed argument supplies no such
checkpoint.

The inspected source resets its supplied seed in both residual sampling and
process generation. That behavior must retain its own method identity and
cannot be described as verified independent streams
\citep[\code{sampleResiduals} and the process-law branches]{chainladder2023snapshot}.
A proposed replacement must specify generator version, stream derivation and
restart behavior, and then check them in execution. Changing a method, law or
request identifier creates a different request fingerprint. A resumed
scheduling request must link to the old checkpoint under the persistence
rules; it cannot silently alter the old request.

The retained snapshot's GPL terms, transitive dependencies and any future
instrumented source remain part of the proposed runtime decision. No package
installation, distribution clearance, crash-durability result or empirical
performance follows from these exchanges. What they establish is narrower
and necessary: a successful payment, a missing value and a failed computation
remain distinguishable when reserving results enter the financial model.


\clearpage
\section{Acquisition costs before the renewal contracts exist}
\label{sec:insurance-acquisition-recoverability}

An insurer may pay acquisition costs that help obtain both current contracts
and their expected renewals. Some of the related future groups do not yet
exist as recognised insurance balances. The retained AASB~17 provisions
therefore require a separate acquisition-asset record for each related group,
unless the permitted immediate-expense election is used. Qualifying costs
are those directly attributable to the insurance portfolio; this is not a
general permission to capitalise selling expenses. A systematic and rational
allocation distinguishes the original group from groups expected to contain
its renewals \citep[paragraphs 28A--28B, B35A and Appendix A]{aasb17}.

Consider a separate, authored example in currency units. Follow only 60 of
qualifying costs already paid and allocated to expected renewals: 40 from
one origination cohort to future group G1, and 20 from another cohort to
future group G2. Other acquisition allocations are outside this example.
Assume that both allocations are supported, neither future group has yet
been recognised, the immediate-expense election is not used, and facts now
indicate possible impairment. Those facts trigger the reporting-date
recoverability assessment under paragraph 28E. An expected renewal remains
a forecast, not a newly recognised premium receivable
\citep[paragraphs 28B, 28E and 35]{aasb17}.

Each future group contains the relevant renewals together with other new
contracts. For G1, suppose the whole group's expected net inflow is 35,
of which the renewal cohort contributes 25. For G2, the corresponding
amounts are 25 and 10. These are conditional measurements under paragraph
32(a), incorporating the required cash-flow, time-value and risk adjustments.
They are neither gross premiums nor CSM. The acquisition cash already paid
does not become a second future outflow. In this additive example, the
other contracts contribute 10 to G1 and 15 to G2.

Paragraph B35D(a) first tests the asset against its related group's expected
net inflow. G1's asset of 40 exceeds 35, so recognise impairment of 5 and
retain 35. G2's asset of 20 is below its whole-group inflow of 25, so this
test produces no loss for G2. Passing that first test does not settle whether
the renewals themselves can recover the costs allocated to them.

The second test addresses that narrower population. Under B35D(b), the
allocated renewal costs exceed the expected renewal inflows by $40-25=15$
for G1 and $20-10=10$ for G2. But 5 of G1's shortfall has already been
recognised. B35D(b)(ii) requires that overlap to be deducted, leaving
additional losses of 10 for G1 and 10 for G2. Total impairment is therefore
$5+10+10=25$, and the remaining assets are 25 and 10. Adding the first loss
to the two unadjusted renewal shortfalls would instead record 30, counting
the same loss of 5 twice \citep[paragraph B35D]{aasb17}.

For these disjoint allocations, write $A_g$ for the original asset, $N_g$ for
the whole-group expected net inflow and $U_g$ for the relevant renewal
inflow. With no prior impairment, the two losses and the resulting balance
are
\[
 I_g=\max(A_g-N_g,0),\qquad
 J_g=\max(A_g-U_g-I_g,0),\qquad
 B_g=A_g-I_g-J_g.
\]
The subtraction of $I_g$ inside the second test preserves the already
recognised loss. This local calculation relies on the identified, nonoverlapping
renewal cohorts and positive inputs. A pool with several overlapping origins
would need an explicit allocation of the relevant costs, inflows and prior
losses before the same reasoning could be applied.

\clearpage
The entries also distinguish impairment from a cash payment. Debit impairment
expense 5 and credit the G1 acquisition asset 5 for the first test. Then debit
impairment expense 20 and credit each group's asset by 10 for the second.
Cash does not move: the acquisition costs were already paid. For a small
balance-sheet slice starting with cash and equity of 100 and no liabilities,
paying the allocated costs leaves cash 40 and acquisition assets 60. After
impairment, cash remains 40, acquisition assets are 35 and equity is 75.
Paragraph 79 includes these acquisition assets within the related issued-contract
portfolio carrying amounts; the separate records used here do not prescribe
a separate line on the published balance sheet
\citep[paragraphs 28E and 79]{aasb17}.

Recoverability can improve. At a later reporting date, suppose the full-group
net inflows are 45 and 30, with renewal contributions of 35 and 15. Assume no
intervening cash movement or recognition of either group. The allocation
assumptions have been reviewed under B35B and remain unchanged; only the
measurement inputs improve. Against the original allocations of 40 and 20,
the remaining shortfalls are now 5 for each group. Paragraph 28F permits
reversal to the extent the impairment conditions have improved. Reverse 10
for G1 and 5 for G2: debit the acquisition assets 15 and credit the impairment
reversal in profit or loss 15. Assets become 35 and 15, totalling 50; cash
stays 40 and equity becomes 90. The improvement restores 15 of the earlier
loss of 25 without creating cash or recognising the expected renewals early
\citep[paragraphs 28F and B35B]{aasb17}.

If the impairment conditions later disappear altogether, with these other
facts unchanged, only the remaining loss of 10 can be reversed; the original
allocated cost of 60 is the ceiling in this example. When related contracts
are actually recognised and the acquisition cash flows enter their group
measurement, the corresponding acquisition asset must be derecognised.
B35C provides for partial incorporation when contracts are added over several
periods. Keeping both the old asset and its incorporated amount would double
count the same cost. The record must therefore preserve its origination cohort,
future group, allocation history, impairment and reversal history, and the
amount already incorporated into recognised contracts
\citep[paragraphs 28C--28F and B35B--B35C]{aasb17}.

These are ordinary reporting-date impairment and reversal entries. They do
not use the separate transition rules for acquisition assets. The historical
edition, actual cost eligibility and supported renewal estimates must be
checked before applying the example to an insurer.

\clearpage
\section{When the acquisition asset enters the insurance group}
\label{sec:insurance-acquisition-incorporation}

Continue from the preceding reversal: cash is 40, the two acquisition assets
are 35 and 15, and equity is 90. The next event is recognition of the related
insurance groups. Assume that coverage begins and the first premiums become
due at this date, with no earlier recognition trigger. All contemplated
contracts are admitted together to their correctly identified groups. No
service has yet elapsed and no claims have been incurred. Paragraphs 25--28
determine recognition; the later cash receipt is a separate event
\citep[paragraphs 25--28]{aasb17}.

Keep the measured group net inflows of 45 and 30 from the previous example.
To expose their components, suppose premiums of 100 and 60 are due immediately,
the present values of future claims and expenses are 50 and 25, and each
group has a nonfinancial risk adjustment of 5. Thus fulfilment cash flows
before premium receipt are $50+5-100=-45$ and $25+5-60=-30$, using positive
amounts for net outflows. These are supplied conditional measurements;
present values do not imply a zero discount rate. Previously paid acquisition
costs are not another future payment. Assume general measurement and no other
assets, liabilities, taxes, reinsurance or intervening measurement changes.

Paragraph 38 requires initial CSM to prevent income or expense from the FCF,
any initial cash flows and the derecognition of the acquisition asset, subject
to its stated exceptions. The paid cost therefore affects initial CSM even
though it is no longer a future cash flow. Let $A$ denote the acquisition
asset's remaining carrying amount, $F$ the FCF and $C$ the CSM. Removing the
old asset without a gain or loss requires the new group balance to replace it.
With positive $L$ denoting a liability, the relation in this case is
\[
 L=F+C=-A,\qquad C=-F-A.
\]
G1 consequently has CSM $45-35=10$ and a net group asset of 35; G2 has CSM
$30-15=15$ and a net group asset of 15. Both margins are positive, so the
onerous-group exception does not arise. There are no other previously
recognised cash-flow assets to which paragraph B123A's revenue-and-expense
exception would apply \citep[paragraphs 28C, 32, 38, 47 and B123A]{aasb17}.

An illustrative recognition entry debits the two group assets by 35 and 15
and credits the corresponding acquisition assets by those amounts. Cash is
still 40 and equity is still 90. The separate acquisition balances are now
zero. The amounts transferred are the carrying values after impairment and
reversal, rather than the original allocations of 40 and 20. The earlier
loss history remains available even though those assets have entered the
group measurements \citep[paragraphs 28C--28D and B35C]{aasb17}.

Now receive the premiums in full, immediately after that measurement. Debit
cash 160 and credit the group carrying balances by 100 and 60. Because the
premiums are no longer future inflows, FCF becomes 55 for G1 and 30 for G2.
CSM stays 10 and 15: no service or time has passed. The groups therefore carry
liabilities of $55+10=65$ and $30+15=45$. Cash has increased to 200, and
$200-65-45=90$ confirms unchanged equity. Paragraph 83 ties revenue to
service, while B123(a)(i) expressly distinguishes premium cash inflows from
the liability reductions that generate revenue. Later acquisition-recovery
revenue and the equal service expense are allocated over time under B125;
neither is created by this instantaneous receipt
\citep[paragraphs 40--41, 83--85 and B123--B125]{aasb17}.

\clearpage
The three balance-sheet states make the transfer visible. Amounts are
currency units; group assets and liabilities below are their net carrying
amounts, not their separate FCF and CSM components.
\begin{center}
\begin{tabular}{lrrr}
\toprule
 & Before recognition & After recognition & After receipt \\
\midrule
Cash & 40 & 40 & 200 \\
Acquisition assets & 50 & 0 & 0 \\
Group assets & 0 & 50 & 0 \\
Group liabilities & 0 & 0 & 110 \\
Equity & 90 & 90 & 90 \\
\bottomrule
\end{tabular}
\end{center}
Receiving 160 has changed liquidity and the insurance balances without
creating profit. Omitting the acquisition assets from the initial CSM
calculation would instead give margins of 45 and 30, overstating total CSM
by 50. Retaining the separate acquisition assets after the correct transfer
would overstate assets and equity by 50.

Even a balanced statement can conceal both mistakes together. If the
implementation leaves acquisition assets of 50 and uses the overstated
margins, its post-receipt liability becomes 160. Cash 200 plus the retained
asset 50 less liability 160 still equals equity 90. That identity passes,
but the asset, liability and CSM are wrong relative to the stated measurement.
Checking the components and the derecognition event is therefore necessary
in addition to checking that the statement balances.

As a boundary case, if no acquisition asset exists, $A=0$ gives $C=-F$ and
a zero group carrying amount before premium receipt, under the same
positive-margin assumptions. A general implementation must also handle
onerous groups, other initial cash-flow balances and contracts added over
time. The present calculation establishes the two-event bridge for the
specified groups, rather than those additional cases or an insurer's actual
valuation.

\clearpage
\section{Acquisition costs and the timing of reinsurance}
\label{sec:insurance-acquisition-treaty}

An acquisition asset can enter a group that is onerous at recognition.
Reinsurance then raises a second question: how much of that initial loss may
be recognised as recoverable, and when? Following the two groups separately
prevents a recoverable loss from becoming either a second asset or a reduction
of insurance revenue.

Consider a separate general-measurement group with opening capital and cash
of 50. Before recognition, the insurer pays acquisition costs of 10, allocated
wholly to that future group. At the prior assessment, expected premium 100,
claims 80 and risk adjustment 5 give expected net inflow 15. The asset of 10
is therefore recoverable under the stated assumptions, leaving cash 40,
acquisition asset 10 and equity 50. No renewal allocation is involved
\citep[paragraphs 28A--28F and B35A--B35D]{aasb17}.

The expected claims increase to 100 when the contracts are recognised;
premium remains 100 and risk adjustment remains 5. The example has no
intervening unrecorded onerous period, discounting, tax, investment component
or OCI. Before premium receipt, fulfilment cash flows are $100+5-100=5$.
Removing the acquisition asset adds 10 to the initial loss, making 15,
with CSM zero. The entry debits loss 15 and credits acquisition asset 10
and remaining-coverage liability 5. Receiving the premium then gives cash
140, liability 105 and equity 35
\citep[paragraphs 28C, 38(c)(i) and 47]{aasb17}.

The loss component of 15 is part of the liability of 105. Acquisition cash
has already left the company: putting it back into future cash outflows
would count the cost twice. Its original asset also disappears when it enters
the group measurement.

A separate held treaty covers half of every claim of this group and costs
55. Assume no commissions, limits, disputes or counterparty shortfall.
Expected recovery is 50; risk transferred is separately valued at 2.5.
That supplied risk adjustment need not equal half the gross adjustment in a
different contract. Ordinary cost CSM is $55-50-2.5=2.5$, so the held
asset after premium payment would be 55
\citep[paragraphs 62--65]{aasb17}.

If the treaty was entered into at or before recognition of the onerous group,
the specified initial loss recovery is $15(0.5)=7.5$. It increases held
cost CSM to 10 and the measured asset to $50+2.5+10=62.5$.
The standard multiplies the recognised loss by the expected claims-recovery
percentage; this does not mean the treaty separately reimburses half the
acquisition invoice. The loss-recovery component of 7.5 is contained in the
asset of 62.5. After paying the treaty premium, cash 85 plus that asset equals
liability 105 plus equity 42.5
\citep[paragraphs 66A--66B and B119C--B119F]{aasb17}.

At the end of the single service interval, release the issued liability of
105 as revenue 90 and loss reversal 15. The ratio $15/105=1/7$ allocates
the expected-claim and risk-adjustment releases to the loss component.
Acquisition-cost presentation adds revenue 10 and an equal service expense
10 under B125. It creates neither cash nor a new acquisition asset.
Recognising and paying actual claims of 100 therefore gives issued-group
service-period profit $100-100-10+15=5$. Including the initial loss of
15, its lifetime loss is 10
\citep[paragraphs 49--52 and B120--B125]{aasb17}.

The held asset releases as premium allocation 55 and loss-recovery reversal
7.5, with incurred recovery and collection of 50 recorded separately.
Its service-period result is $50-55-7.5=-12.5$; the initial recovery
of 7.5 makes its lifetime result $-5$. The company's final cash and equity
are both 35:
\[
 50+100-10-55-100+50=35.
\]
All issued and held balances are zero. Total operating cash is $-15$,
profit is $-15$, and OCI is zero. Premium allocation remains separate from
issued revenue \citep[paragraph 86]{aasb17}.

If the otherwise identical treaty is entered into only after the issued
group is recognised, B119C prevents recovery of that original loss.
Its initial held asset is 55 and its cost CSM is 2.5. Equity is initially
35 and remains 35 after the combined subsequent service result of zero.
The timely treaty instead gives initial equity 42.5 and a subsequent loss
of 7.5. Equal final cash does not make those recognition histories
interchangeable.

The transfer therefore requires the acquisition allocation, recoverability
history and unconsumed asset version. The held measurement additionally
requires treaty entry time, covered-contract mapping, claims recoverability
and its own risk adjustment. Exact replay does not transfer the asset again;
an unresolved treaty date does not default to timely recognition. Partial
coverage or changing estimates require their stated allocations before the
same procedure can be applied.



\section{Reading acquisition costs in an actual report}
\label{sec:munich-new-contracts}

Section~\ref{sec:insurance-acquisition-incorporation} supplied the acquisition assets and the net inflows
for two identified groups. A published insurer's accounts reveal a different
level of detail. Munich Re's 2023 report shows contracts recognised for the
first time during the year, separating concluded from acquired contracts and
nononerous from onerous business. Note 27 covers segments predominantly using
general measurement or the variable fee approach; it omits two predominantly
PAA segments. These are segment disclosures, not a list of individual groups
\citep[printed pp. 252--254]{munichre2023annual}.

Consider the ERGO Life and Health Germany segment, identified in the paragraph
preceding its table. Its newly recognised contracts combine GMM and VFA
measurements. The following two columns reproduce its concluded-contract
amounts in millions of euros. Preserve the report's signs: inflows are positive,
outflows and the margin are negative. Its onerous CSM cell is blank; a dash
below preserves that blank instead of assigning a reported zero
\citep[printed pp. 252--253]{munichre2023annual}.

\begin{center}
\begin{tabular}{p{0.58\textwidth}rr}
\toprule
 & Nononerous & Onerous \\
\midrule
Present value of future inflows & 2,842 & 17 \\
Future claims, expenses and investment components & $-2,438$ & $-12$ \\
Future acquisition costs & $-243$ & $-25$ \\
Risk adjustment & $-8$ & 0 \\
Contractual service margin & $-154$ & --- \\
Reported total & 0 & $-21$ \\
\bottomrule
\end{tabular}
\end{center}

The nononerous amounts illustrate how expected future profit is deferred in
CSM. The onerous column retains a loss-bearing population separately. Munich
Re describes recognising an expected initial loss in profit or loss and
keeping onerous contracts separate from profitable ones. Combining these
columns into one net profit would erase that distinction
\citep[printed pp. 165--166]{munichre2023annual}.

Adding the displayed nononerous numbers gives $-1$, although the table reports
zero. The onerous cash-flow and risk-adjustment rows give $17-12-25-0=-20$,
against a reported total of $-21$. The report states that amounts are rounded
to millions and that totals may consequently differ. A possible set of
unreported nononerous values is $2{,}842.3-2{,}437.8-242.8-7.8-153.9=0$;
each rounds to its displayed amount. Similarly,
$16.6-12.1-25.2-0.3=-21$ is consistent with the onerous column if no CSM is
assigned. Together these hypothetical values also reproduce the table's
outflow subtotals and aggregate rows after rounding. They demonstrate
consistency under nearest-million rounding; they do not recover the issuer's
unreported measurements or establish its exact rounding convention. Retain
the original figures and blank cells \citep[printed pp. 155, 253]{munichre2023annual}.

\clearpage
The acquisition-cost row appears within \emph{future cash outflows}. Its
reported total of $243+25=268$ therefore cannot simply be entered as the
previous example's asset $A$, which arose from costs already paid before the
related groups were recognised. Nor do the table's acquired-contract columns
describe how sales acquisition costs were paid. Those columns concern a
different distinction between contract populations. Their displayed zeros
retain the same rounding qualification. The selected disclosure does not
identify separate prior acquisition assets, their impairment history or their
allocation to individual groups \citep[printed pp. 252--253]{munichre2023annual}.

A third quantity appears in the segment's service revenue. For 2023, Munich
Re reports €202 million as the portion of premium relating to recovery
of acquisition costs. Its policy explains that acquisition costs enter
insurance service expense in the same amount as the recovery included in
revenue, allocated over the coverage period. This is a period service amount;
the €268 million above is an estimate of future acquisition outflows for
newly recognised contracts. Subtracting one from the other would not produce
a remaining acquisition asset. The relevant populations, dates and movements
must first be linked \citep[printed pp. 177, 253, 264]{munichre2023annual}.

The segment's liability reconciliation makes that caution concrete: it shows
€637 million in the service-expense row labelled ``Recovery of acquisition
costs'', and €978 million of acquisition costs within cash flows. Thus the
€202 million recovery in revenue does not describe the entire expense row.
The difference of €435 million is an arithmetic residual; assigning its
cause requires the relevant populations and expense treatments. Neither that
residual nor cash less expense identifies a separate acquisition asset
\citep[printed pp. 230--231]{munichre2023annual}.

The same timing distinction applies to profit. The report gives
€819 million of CSM for services provided during 2023 in this segment.
That amount concerns service during the year, whereas €154 million is
the margin attached to the newly recognised contracts in the selected table.
It is not a release from that new-business amount alone. A reconciliation
needs the opening CSM and every intervening movement. The segment's total
insurance revenue also combines measurement approaches: €7,966 million
under GMM/VFA and €1,976 million under PAA make the reported
€9,942 million \citep[printed pp. 253, 264]{munichre2023annual}.

Even within one insurer, an acquisition-cost policy needs its applicability
conditions. Munich Re reports using the PAA option to expense acquisition
costs directly for a substantial part of its property-casualty primary
insurance business. The earlier general-measurement asset example therefore
requires the relevant treatment and group facts before it can be applied to
a real portfolio \citep[printed p. 166]{munichre2023annual}. These are readings
of the retained 2023 English translation, whose German original is binding.
They identify reported aggregates and necessary further facts; the underlying
group histories and valuations remain unresolved.

\clearpage
\section{Following the margin through the year}
\label{sec:munich-csm-bridge}

Why can the segment release €819 million of CSM during 2023 when newly
recognised contracts add only €154 million? The opening margin carries profit
from earlier business. Munich Re's measurement-component reconciliation for
ERGO Life and Health Germany begins with CSM of $-10{,}449$ million euros
and ends with $-9{,}569$ million. Its signs present the margin as a negative
balance: new business makes it more negative, while service release moves
it towards zero. Keep that convention while reading the movements
\citep[printed p. 246]{munichre2023annual}.

The nonzero CSM movements are service provided during the year, contracts
initially recognised, changes in estimates that adjust CSM, and derecognition
or changes in the consolidated group. In that order, the displayed amounts
give the complete bridge
\[
 -10{,}449+819-154+214+1=-9{,}569.
\]
In positive margin amounts, opening CSM of 10,449 receives 154 from new
business, releases 819 for service, and falls by another 214 and 1 through
the other two movements. The net reduction is 880. The table's CSM service
subtotal of 879 already combines $819-154+214$; adding that subtotal to the
individual movements would count them twice. The estimate movement here is
the CSM column's 214, not a claim that every change in the underlying
estimates is absorbed by CSM \citep[printed p. 246]{munichre2023annual}.

The numbers connect three distinct disclosures. New-business CSM of 154
matches the initial-recognition table; the 819 release matches the segment's
service revenue; and closing CSM of 9,569 is the balance used in the
issuer's schedule of expected future releases. The latter match verifies
which balance that schedule starts from, rather than the accuracy of the
future releases. The comparative-year column also reconciles:
$-9{,}412+948-221-1{,}764=-10{,}449$, which becomes the 2023 opening
amount \citep[printed pp. 247, 253--254 and 264]{munichre2023annual}.

The larger liability totals require a different comparison. The measurement
reconciliation includes the full GMM/VFA balances and the incurred-claims
liabilities of PAA contracts. It excludes PAA remaining coverage. It would
therefore be wrong to describe the whole table as GMM/VFA only. The footnote
continued after its comparative table explicitly gives the omitted PAA
remaining-coverage balance: $-62$ million euros in the comparative period
and $-1$ million at the end of 2023 \citep[printed pp. 241 and 247]{munichre2023annual}.

This scope difference accounts for the displayed opening and closing
differences. The complete liability reconciliation starts at $-109{,}556$,
against $-109{,}494$ in the measurement table: the difference is $-62$.
Their closing amounts are $-112{,}808$ and $-112{,}807$, a difference of
$-1$. Thus an apparently small missing balance is a disclosed excluded
population, rather than permission to omit it or absorb it into rounding
\citep[printed pp. 230--231 and 246--247]{munichre2023annual}.

\clearpage
The cash rows preserve the same distinction. Premium receipts change the
reported liability balance by $-10{,}406$ in the full table and $-8{,}471$
in the measurement table, leaving $-1{,}935$ attributable to the omitted
PAA remaining-coverage portion at displayed precision. Acquisition cash
movements of 978 and 502 leave 476 for that portion. These signs describe
the liability reconciliation; they are not the debits and credits of a
cash account. Nor are these differences all PAA cash flows, because PAA
incurred claims are included in both tables
\citep[printed pp. 231, 241 and 246]{munichre2023annual}.

The 476 cash difference still does not explain the earlier acquisition
expense difference of $637-202=435$. The two quantities differ by 41, but
calling that amount a deferred acquisition asset would again invent a
missing link. A cash payment, an expense and recovery included in service
revenue can concern different dates and contract populations. The disclosed
expensing election makes the treatment of PAA costs a relevant question;
the selected tables do not establish that it alone explains 435. Preserve
that unresolved attribution alongside the now-reconciled CSM and the
identified PAA scope difference
\citep[printed pp. 166, 230--231, 246 and 264]{munichre2023annual}.

All these comparisons use the report's rounded euro millions. Exact addition
of the CSM column does not establish unrounded measurement equality or
exact conservation of every other column. The public reconciliation supplies
movements and their reporting populations, while the group assignments,
valuation inputs and detailed acquisition-expense allocation remain separate
facts to obtain.

\clearpage
\section{Adoption dates and the restatement of 2022}
\label{sec:insurance-adoption}

A standard's effective date and the economic date of a reported balance need
not coincide. An insurer adopting new rules in 2023 may report a revised
balance for December 2022 while retaining the earlier report as a record of
what it originally disclosed. The system must preserve both versions and
explain which one a calculation uses.

The Ministry of Finance's explanation dated 12 June 2025 describes the original
ASBE~25 timetable. It specifies application from 1 January 2023 for enterprises
listed both domestically and overseas, and for overseas-listed enterprises
preparing financial statements under IFRS or ASBE. Other enterprises applying
ASBE were scheduled to apply the revised standard from 1 January 2026.
Its account of the
2025 notice says existing adopters, including enterprises that had adopted
early under the 2020 notice, must continue applying it. Insurers facing genuine
difficulties and seeking deferral must submit written reasons to the Ministry's
Accounting Department and the relevant institutional supervision departments
of the National Financial Regulatory Administration. That statement does not
establish an automatic deferral merely because preparation is difficult
\citep[answers on background and implementation arrangements]{mofinsurance2025qa}.

The same explanation limits the new simplification choices to unlisted
enterprises; matters outside those choices remain subject to ASBE~25. It names
grouping, measurement, presentation and transition as affected areas, including
comparative information and the transition date. These are source-supported
descriptions of the Ministry's explanation. The operative 2020 and 2025 notices
have not yet been directly verified in this collection, so the explanation
does not supply the detailed conditions needed to implement those choices
\citep[answer on the relationship between the notice and the standard]{mofinsurance2025qa}.

Ping An provides a concrete adoption history. Its 2022 ASBE annual report
states that the group would adopt the revised insurance standard for accounting
years beginning on 1 January 2023. Its 2023 report confirms adoption on that
date and restatement of comparative information
\citep[Note XIV.2, pp.~315--316]{pingan2022asbe}
\citep[Note IV.34, p.~192]{pingan2023asbe}. The following balances all refer
to the consolidated group at 31 December 2022, in RMB millions. The first
column comes from the 2022 report; the second is the restated comparative in
the 2023 report. Changes are calculated from those columns; the asset,
liability and parent-equity changes also match the three rows of Note IV.34's
reported bridge \citep[pp.~146--148]{pingan2022asbe}
\citep[pp.~141--143 and 192]{pingan2023asbe}.

\begin{center}
\small
\begin{tabularx}{\linewidth}{@{}Xrrr@{}}
\toprule
Balance & Original & Restated & Change\\
\midrule
Assets & 11,137,168 & 11,009,940 & $-127,228$\\
Liabilities & 9,961,870 & 9,823,944 & $-137,926$\\
Total equity & 1,175,298 & 1,185,996 & $10,698$\\
Parent shareholders' equity & 858,675 & 869,191 & $10,516$\\
Non-controlling interests & 316,623 & 316,805 & $182$\\
\bottomrule
\end{tabularx}
\end{center}

The distinction between total equity and equity attributable to the parent
matters. Assets less liabilities rise by
$-127,228-(-137,926)=10,698$. Of that increase, 10,516 belongs to parent
shareholders and 182 to non-controlling interests. Treating the parent-equity
row as total equity would leave an unexplained difference of 182. These are
changes between reported accounting bases for the same date, not cash received
in a later period.

The transition also changes how later growth is calculated. Ping An reports
assets of 11,583,417 at the end of 2023. Subtracting its restated 2022
comparative gives a reported-balance increase of
$11,583,417-11,009,940=573,477$. Subtracting the original 2022 amount instead
gives 446,249, mixing the old and new bases and understating that comparable
movement by 127,228 \citep[p.~141]{pingan2023asbe}. The first difference is
a movement in reported asset balances; it is not a cash-flow measure or an
explanation of every economic cause of growth.

An original-source query for assets at 31 December 2022, explicitly selecting
the 2022 report, therefore returns 11,137,168. A query selecting the 2023
report's restated comparison for that date returns 11,009,940. Each assertion
retains its report identity, consolidated scope, RMB-million scale, accounting
basis and economic date. A link records that the later assertion restates the
comparison; it does not delete the original. A request for a comparable annual
movement must reject the mixed-basis pair unless an explicit reconciliation
has first aligned it. Because the basis and disclosure version differ, these
two amounts are not a same-context source conflict.

The information date remains a separate question. The 2023 report says its
board authorised the statements on 21 March 2024
\citep[Note I, p.~154]{pingan2023asbe}. Neither that authorisation nor the
31 December reporting date establishes when the exact retained document became
public. The corpus has archive-date leads but no verified publication date or
historical byte vintage for either selected report. A query for information
available to investors at 31 December 2023 must therefore retain an unresolved
availability result for this source selection; it cannot silently import the
later restated balance. The corpus's October 2026 revalidation timestamps
describe local evidence handling, not investor knowledge in 2023.

Finally, Note IV.34 reports financial-asset reclassification, including use of
the classification overlay for relevant assets derecognised during 2022. It
also states that some insurance groups used modified-retrospective or
fair-value transition because retrospective application was impracticable.
The headline bridge consequently cannot be reconstructed by applying one
percentage to every insurance liability. The example establishes how to retain
and compare disclosed histories; reconstructing the transition itself still
requires the affected groups, asset classifications and underlying measurement
inputs \citep[Note IV.34, p.~192]{pingan2023asbe}.

\clearpage
\section{Reconstructing balances at transition}
\label{sec:insurance-transition}

A disclosed restatement tells us what balances an insurer reported under the
new basis. It does not reveal the historical contract groups, assumptions or
services needed to reproduce their opening measurement. The three transition
methods address different limits on that reconstruction. The retained
December 2022 AASB~17 compilation provides a useful, explicitly scoped account
of those methods: its stated applicability covers annual periods beginning
from 1 January 2023 to before 1 July 2026. These provisions explain a design
requirement; they do not establish the rules governing Ping An's ASBE groups
\citep[Appendix C]{aasb17}.

First distinguish the dates. Initial application begins the annual period in
which the entity first applies the standard. Transition normally begins the
immediately preceding annual period. Thus a calendar-year entity first applying
the standard in 2023, with only 2022 presented as an adjusted comparative,
starts its transition calculation on 1 January 2022. If it elects to present
earlier adjusted comparatives, paragraph C25 instead takes transition back to
the beginning of the earliest such period. Neither date establishes when the
resulting report became public \citep[paragraphs C2 and C25--C27]{aasb17}.

Full retrospective application identifies and measures each group as if the
standard had always applied. Acquisition-cash-flow assets are considered
separately, balances that would not exist under the new rules are removed, and
the resulting net difference is recognised in equity. This route needs the
historical information required to reconstruct grouping, cash flows,
measurement assumptions and service already provided. It is required unless
application is impracticable or the specific paragraph C5A exception applies
\citep[paragraphs C3--C5]{aasb17}.

Impracticability requires more than a missing historical workbook. Under
AASB~108, the entity must have made every reasonable effort to apply the
requirement. Relevant difficulties include effects that cannot be determined,
assumptions that would have to be made about management's past intentions,
and significant estimates for which the necessary contemporaneous evidence
cannot be objectively distinguished from other information. Those conditions
need to be established for the reconstruction at issue
\citep[paragraph~5]{aasb108comp2022}.

An estimate prepared later can nevertheless describe an earlier period. The
question is whether it reflects the circumstances that existed then, using
evidence that would have been available when the earlier financial statements
were authorised for issue. A newly assembled estimate can meet that test;
replacing an earlier expectation with knowledge of what subsequently happened
uses hindsight. The date of preparation alone therefore does not settle the
question. Nor does authorisation establish when an outside investor could
obtain the information: that requires the separate public-availability record.
The insurance standard's specific transitional provisions govern their stated
cases, including permitted information under the modified and fair-value
routes described below
\citep[paragraphs~7, 19(a) and 50--53]{aasb108comp2022}.

When full retrospection is impracticable for a group, the alternatives are
modified retrospection and fair-value transition. Modified retrospection aims
for the closest retrospective result possible using reasonable and supportable
information available without undue cost or effort. Each specified modification
is allowed only to the extent that the necessary retrospective information is
missing. If the information needed even for this modified method cannot be
obtained, paragraph C6(a) requires the fair-value approach. Missing history
therefore does not license an unrestricted choice of present assumptions
\citep[paragraphs C5--C8]{aasb17}.

For example, the modified route can reconstruct earlier expected cash flows
from estimates at transition, or an earlier determinable date, adjusted for
known intervening cash flows, including those of contracts that have expired.
Its discount-rate approximations require the specified yield-curve evidence;
its risk adjustment and allocation of past service require their own records.
The source also prescribes nil amounts for particular acquisition-cash-flow
items when the necessary information is unavailable. That narrow instruction
cannot be generalized into a rule that every missing amount is zero
\citep[paragraphs C12--C16]{aasb17}.

Fair-value transition takes a different starting point. Paragraph C20 determines
the CSM or loss component from the difference between the fair value of the
insurance group and its fulfilment cash flows, both measured at transition.
The insurance group's fair value is distinct from the value of its underlying
investment pool. Replacing one with the other would compute a different
quantity. The calculation needs a supported group valuation as well as the
cash-flow, discount and risk-adjustment inputs; C20 expressly excludes the
demand-feature rule in AASB~13 paragraph~47. Its permitted information-date,
grouping and accumulated-OCI choices also have specified conditions
\citep[paragraphs C20--C24B]{aasb17}.

To interpret that group valuation, consider what would be paid to transfer
the liability in an orderly transaction between market participants at the
measurement date. The obligation remains outstanding, and the transferee must
fulfil it; the thought experiment is not an immediate settlement with the
policyholder. AASB~13 therefore calls for assumptions that market participants
would use, in the accessible principal market or, in its absence, the most
advantageous market. The unit being valued still comes from C20: the insurance
group \citep[paragraphs~2--3, 9, 13--24 and 34]{aasb13comp2021}.

A quoted price is not required in every case. The valuation should maximise
relevant observable inputs, while supported unobservable inputs can supply
the participant assumptions for which market observations are unavailable.
The insurer's own data can be a starting point, but require adjustment when
reasonably available evidence indicates different participant assumptions or
entity-specific features unavailable to other participants. A technique must
also fit the circumstances and available data. These requirements make the
source, date and interpretation of each valuation input consequential
\citep[paragraphs~61--67 and 87--89]{aasb13comp2021}.

There is a further reason to retain separate fair-value and FCF calculations.
AASB~13 includes non-performance risk, including the issuer's own credit
risk, in liability fair value and assumes that risk is unchanged by the
hypothetical transfer. AASB~17 expressly excludes the issuer's own
non-performance risk from fulfilment cash flows for issued insurance contracts.
Even when the two measurements draw on some of the same cash-flow records,
their treatment of this risk differs. Neither a single imported FCF balance
nor the value of the underlying investments supplies both C20 inputs
\citep[paragraphs~42--44]{aasb13comp2021}
\citep[paragraph~31]{aasb17}.\footnote{The comparison uses AASB~108's
31 December 2022 compilation, applying to periods beginning from 1 January
2023, and AASB~13's 30 June 2021 compilation, applying from 1 July 2021.
These are retained historical editions; later amendments and the actual
entity's applicable requirements need their own check.}

\clearpage
There is one important qualification to the ordinary method choice. Under
C5A, a group with direct participation features may elect fair-value transition
even when retrospective application is possible, if the entity chooses the
B115 risk-mitigation option prospectively from transition and used qualifying
mitigation before that date. The applicable instrument, qualification and
designation conditions still matter. Treating every fair-value election as
proof of impracticability would therefore be wrong
\citep[paragraphs C3(b), C5A and B115--B118]{aasb17}.

The reconstruction must consequently retain the method and its basis for each
group, each modification used, the information date and evidence supporting
it, the separately measured FCF, CSM or loss component, relevant OCI and
acquisition balances, and the opening-equity reconciliation. Some transition
rules expressly permit transition-date information for specified assessments.
That permission concerns the accounting restatement; it does not make a later
report admissible in a forecast made before its publication. Detailed
impracticability judgments under AASB~108 and valuation under AASB~13, together
with actual group facts, remain necessary before this method description can
produce a supported opening measurement.


\clearpage
\section{From transition measurements to opening equity}
\label{sec:insurance-transition-balances}

The fair-value approach produces a margin or a loss component, but those two
outcomes enter the liability differently. Consider two distinct groups of
issued contracts for which the grouping and eligibility for this transition
method have already been established. Use the historical provisions just
described and assume that both groups contain only remaining coverage. All
amounts below are authored currency units at the same transition date. There
are no incurred claims, reinsurance, acquisition balances, tax, accumulated
OCI or other transition adjustments in this example. These simplifications
isolate the effect of the two measurements; they do not supply actual insurer
valuations or evidence of eligibility.

For group A, suppose the measured fulfilment cash flows are 90 and the
separately measured fair value is 110. The FCF amount already includes the
discounting and nonfinancial risk adjustment required by the measurement
rules. Paragraph C20 makes the difference, $110-90=20$, the CSM. This is
unearned profit associated with future services. Keeping it within the
liability gives $90+20=110$; recording it as immediate profit would contradict
that role \citep[paragraphs 32--38, 40(a) and C20]{aasb17}.

Now reverse the ordering of the inputs. For group B, let FCF be 120 and fair
value 100. The difference identifies a loss component of 20, while CSM is
zero. The liability is 120. Paragraphs 49--50 describe the loss component as
part of the liability for remaining coverage, with its own subsequent
allocation and reversal rules. Adding another 20 to the FCF balance would
count that component twice. Replacing the liability by fair value, 100,
would instead understate the required FCF balance
\citep[paragraphs 40, 47, 49--50 and C20]{aasb17}.

Write $F$ for measured future-service FCF, $V$ for the group's fair value,
$C$ for CSM and $K$ for the nonnegative loss component. Positive amounts
denote liabilities. The two sign cases give the following local derivation:
\[
 C=\max(V-F,0),\qquad K=\max(F-V,0),\qquad
 L=F+C=\max(F,V).
\]
Here $L$ is the recognised remaining-coverage liability. The loss component
$K$ is tracked within $L$ and is not added to it. Thus $C-K=V-F$, and at
$V=F$ both components are zero. These relations apply to the stated
transition case; later changes in FCF, service provision and finance effects
follow the subsequent-measurement rules.

The numerical comparison makes the distinction visible:
\begin{center}
\begin{tabular}{lrr}
\toprule
Currency units at transition & Group A & Group B \\
\midrule
Fulfilment cash flows & 90 & 120 \\
Fair value of the group & 110 & 100 \\
CSM & 20 & 0 \\
Loss component within the liability & 0 & 20 \\
Recognised liability & 110 & 120 \\
\bottomrule
\end{tabular}
\end{center}
The recognised liability equals fair value for A and exceeds it for B.
Calling the method fair-value transition therefore does not mean that every
resulting liability is carried at fair value.

\clearpage
To see the opening-equity effect, suppose the old carrying liabilities are
105 for A and 115 for B. Assets of 500 remain unchanged, and these are the
only liabilities. Before transition, equity is $500-105-115=280$. Afterwards,
it is $500-110-120=270$. Each group's liability increases by 5, so total
opening equity falls by 10. Paragraph C5 replaces the retrospective
measurement requirement in C4(a); the removal of incompatible old balances
and the net adjustment to equity in C4(b)--(c) still apply. The ordinary
initial-recognition loss rule does not turn this transition adjustment into
current-period operating expense \citep[paragraphs C4--C5]{aasb17}.

Illustrative component entries show the same reconciliation. For A, debit
the old liability 105 and opening equity 5, then credit the FCF liability 90
and the CSM liability 20. For B, debit the old liability 115 and opening
equity 5, then credit the FCF liability 120. Record B's loss component of 20
within that 120 for its subsequent allocation; it produces no additional
credit. Neither entry changes cash. An actual transition would also need the
acquisition, reinsurance, tax and OCI effects excluded here.

The order of calculation matters as much as the signs. If an implementation
first combined the two already distinct groups, it would find total FCF of
210 and total fair value of 210. Applying the sign rule to those totals would
produce neither CSM nor a loss component and would report a liability of 210.
The correct group-by-group calculation gives 230. Pooling would therefore
understate liability and overstate equity by 20. The programme must retain
the admitted group boundary while calculating CSM and loss components, then
aggregate the resulting balances. Whether different grouping is itself
permitted requires the separate grouping rules; an aggregation shortcut
cannot decide that accounting question.

The example supplies the arithmetic once the group's two measurements and
old balances are known. A real transition disclosure may reveal only some
of those quantities. The later Munich Re case shows both what can be read from an
insurer's report and what still requires contract and valuation records.

\clearpage
\section{The recoverable part of a transition loss}
\label{sec:reinsurance-transition-recovery}

Group B's loss component of 20 describes the issued contracts. If the insurer
also holds reinsurance, how much of that loss is expected to be recovered?
Extend the preceding example by adding a separately recognised, qualifying
group of reinsurance contracts held. Assume it covers all contracts in B and
only B, with remaining coverage and no incurred claims, and uses general
measurement. This changes the earlier no-reinsurance assumption. The issued
obligation to policyholders remains 120, with its loss component of 20;
purchasing reinsurance does not extinguish it
\citep[paragraphs 60--64 and 75]{aasb17}.

The transition calculation has a specific timing exception. Ordinarily,
paragraph B119C permits the recovery-income treatment in paragraph 66A only
if the reinsurance contract was entered into before or when the onerous
underlying contracts were recognised. Suppose the underlying contracts'
recognition date, applying these rules, is 1 July 2020, the treaty was entered
into on 1 October 2020 and
both coverages remain in force at transition on 1 January 2022. The treaty
does not meet that ordinary condition for the original underlying recognition.
Paragraph C20A nevertheless expressly removes B119C's condition when
determining the loss-recovery component at transition. The exception does
not remove the treaty's other recognition, grouping, boundary or measurement
requirements \citep[paragraphs 61--66B, B119C and C20A]{aasb17}.

Let $q$ be the percentage of claims on the underlying contracts that the
insurer expects to recover from this reinsurance group, expressed as a
fraction. Assume $q=0.60$ for this conditional illustration. A nominal treaty
quota alone does not establish that expectation: coverage, limits, exclusions
and collectability matter. The full reinsurance measurement also includes
issuer non-performance risk and the nonfinancial risk transferred, under
paragraphs 63--64. With $K=20$ still denoting B's loss component, C20A gives
the reinsurance loss-recovery component $R$:
\[
 R=Kq=20\times 0.60=12.
\]
This 12 belongs to the reinsurance asset for remaining coverage. It is not
the issued liability, the amount of expected claims, or a valuation of the
whole reinsurance asset. If only some contracts in the underlying group were
covered, the required allocation of the losses would also have to be
established; the all-covered assumption avoids that extra calculation here
\citep[paragraphs 63--64, 66B, B119E and C20A]{aasb17}.

The complete asset still needs the treaty cash flows, discounting, transferred
risk and net cost or gain of purchasing reinsurance. The reinsurance CSM
reflects that net cost or gain; held reinsurance is not classified as onerous
under paragraph 68. Applying the issued-group CSM and loss formulas to the
held contract would therefore be wrong. Presentation likewise keeps the
issued and held-reinsurance portfolios separate
\citep[paragraphs 65--68 and 78--79]{aasb17}.

\clearpage
To extend the opening balance sheet, supply the missing total explicitly.
Assume a separate complete measurement gives a reinsurance asset of 50,
already including the loss-recovery component of 12, against an old asset of
45. These are additional conditional inputs; the multiplication above does
not derive them. Other assets remain 500 and the two issued liabilities still
increase from 220 to 230. Total assets therefore rise from $500+45=545$ to
$500+50=550$, while equity moves from $545-220=325$ to $550-230=320$.
The asset increase adds 5 to equity against the issued-liability increase
of 10, giving a net reduction of 5. An illustrative reinsurance entry debits
the new asset 50, credits the old asset 45 and credits opening equity 5.
Adding another asset of 12 would count a component of the supplied 50 twice
and overstate equity by 12. The net transition change belongs in equity
under C4(b)--(c), rather than being inferred as current recovery income
from the ordinary recognition rule \citep[paragraphs C4--C5 and C20A]{aasb17}.

The component also imposes a useful later check. If B's loss component falls
to 14 and the recoverable portion remains 60\%, paragraph B119F limits the
reinsurance loss-recovery component to $14\times0.60=8.4$. Leaving it at 12
would violate that bound. The bound alone neither selects a unique new
component nor determines the full reinsurance asset or period's profit;
the actual changes and allocation rules still matter. Similarly, $q=0$
gives no loss-recovery component and $q=1$ gives 20, without erasing the
issued obligation or determining the entire asset
\citep[paragraphs 66--67 and B119F]{aasb17}.

\clearpage
\hypertarget{munich-re-transition-disclosure}{%
\section{What an insurer's transition disclosure reveals}\label{method-insurance-munich-transition}}

Munich Re's 2023 report shows how much a public disclosure can resolve.
Its consolidated statements apply IFRS as adopted by the European Union,
together with the stated German commercial-law requirements. The first
application of IFRS~17 is 1 January 2023, with transition at 1 January 2022
and restated insurance comparatives. Financial-instrument comparatives
generally remain under IAS~39, subject to the disclosed classification
overlay. A single year label therefore does not imply a single measurement
basis for every line in the comparative balance sheet
\citep[pp.~155 and 183--184]{munichre2023annual}.\footnote{The retained
English report states that it is an informational translation and the German
original is binding (p.~1). The discussion attributes the disclosed policies
and amounts; the historical AASB texts above are comparison sources.}

The report identifies full retrospection for certain property-casualty
contracts and parts of life reinsurance. Where historical information is
insufficient in life reinsurance, it mostly uses fair-value transition; it
also uses that method for specified property-casualty reinsurance groups and
German primary business. Modified retrospection and fair value together
account for over 90\% of net insurance liabilities at transition. That is a
statement about a liability denominator at a particular date. It supplies
neither a count of contracts nor the assignment and supporting history of
every group \citep[p.~183]{munichre2023annual}.

The opening-equity bridge makes the scale of the different effects visible.
The reported margin component is much larger than the net equity reduction
because the other measurement changes also enter the reconciliation. The
following amounts reproduce the report's signs and billion-euro scale
\citep[p.~184]{munichre2023annual}.

\begin{center}
\begin{tabular}{lr}
\toprule
Opening-equity reconciliation at 1 January 2022 & EUR billions \\
\midrule
Equity under IFRS 4 & 30.9 \\
Modified measurement basis & 25.3 \\
Contractual service margin & $-22.3$ \\
Risk adjustment for non-financial risk & $-5.6$ \\
Loss component & $-1.4$ \\
Deferred taxes & 1.5 \\
Reported change in equity & $-2.5$ \\
Equity under IFRS 17 & 28.5 \\
\bottomrule
\end{tabular}
\end{center}

The five displayed change components sum to $-2.5$ billion, whereas the
displayed endpoints differ by $-2.4$ billion. The report warns that rounding
can cause such deviations; these figures are compatible under rounding to
the displayed precision. The representation should retain both the reported
endpoints and the bridge, with their precision, instead of changing one to
force an exact identity \citep[pp.~1, 155 and 184]{munichre2023annual}.

\clearpage
The separate margin tables also clarify which contracts are being followed.
At 1 January 2022 the reported issued-contract CSM is $-22{,}688$ million
and the held-reinsurance net cost/gain is $368$ million. Adding those reported
totals gives $-22{,}320$ million, consistent with the rounded $-22.3$ billion
margin component above. The issued-contract transition table follows groups
in force at transition, and its CSM amounts concern the general measurement
model and VFA; the associated revenue also includes PAA contracts. Its
later balances therefore cannot be substituted for the whole subsequent
portfolio or used as if CSM and revenue had identical populations
\citep[pp.~210--211 and 254--255]{munichre2023annual}.

Finally, disclosure of a method does not supply all its inputs. Munich Re
reports a 6\% cost-of-capital rate for its non-financial risk adjustment,
applied to projected capital requirements with diversification and discounting.
Its equivalent one-year confidence level is about 90\% at transition, compared
with 88\% for 2023 and 87\% for the comparative period. These are dated
properties of the insurer's calculation, not universal IFRS parameters.
The rate also cannot be imported into the separate fair-value calculation
merely because that discussion mentions cost of capital. Reconstructing an
individual group still requires its cash flows, capital projections, curves,
allocations and supported transition-method choice
\citep[pp.~165, 183 and 185]{munichre2023annual}.

\clearpage
\hypertarget{interfaces-exclusions-and-acceptance-cases}{%
\section{Interfaces, exclusions and acceptance cases}\label{method-insurance-methods-interfaces-exclusions-and-acceptance-cases}}

Insurance events are group\_recognize, premium\_settle, service\_release, claim\_incur, claim\_settle, assumption\_change, onerous\_change and treaty\_change, each bound through the event-basis, authorization and commit contracts to the exact group/version, evidence and information view. Journals and component roll-forwards must reconcile independently; a component reallocation is not always a net liability change. Invalid grouping, missing rate/law, unsupported option valuation, ambiguous contract boundary or unknown service units produces the corresponding unresolved/outside-scope result. Failed numerical paths retain their coverage/mass, without successful-only normalization. Modified contracts need the applicable derecognition analysis; no new ID can erase unfavorable old-group obligations.

Initial transition requires an actual method choice and eligibility for full retrospective, modified-retrospective or fair-value transition under AASB~17 Appendix C, with acquisition, grouping, past service and historical-rate records. This collection does not claim an executed transition or a general VFA valuation. The nonzero-rate general case and onerous allocation through expiry above are finite specification examples. Required later acceptance still includes a priced guarantee, acquisition/reinsurance interaction, transition and actuarial backtesting. These retained obligations prevent the two-year illustration from being described as full insurance-domain readiness.


